\subsection{超立方体粗粒化的规范体积与边界切割：体--界字典的离散等周接口}\label{subsec:fold-hypercube-bulk-boundary-gauge-volume}
本小节在与定理 \ref{thm:fiber-gauge-volume-chain-factorization} 一致的纤维规范群口径下，讨论超立方体微态空间 \(\Omega_n=\{0,1\}^n\) 上的任意满射粗粒化 \(f:\Omega_n\twoheadrightarrow X\)。
核心目标是将纤维规范体积密度与界面切割面积联系起来，并以边界邻接多图的回路圆维（自由秩）给出等价表述。

\begin{definition}[超立方体粗粒化的界面面积与规范体积密度]\label{def:hypercube-coarsegraining-area-gauge-volume}
记 \(Q_n=(\Omega_n,E_n)\) 为 \(n\)-维超立方体图（相邻当且仅当 Hamming 距离为 \(1\)），并取满射 \(f:\Omega_n\twoheadrightarrow X\)。
对 \(x\in X\) 记纤维 \(S_x:=f^{-1}(x)\) 与纤维大小 \(d_x:=\card{S_x}\)。
定义纤维规范群
\[
\mathcal G(f):=\mathcal G(f:\Omega_n\twoheadrightarrow X)
=\prod_{x\in X}\mathrm{Sym}(S_x),
\qquad
\card{\mathcal G(f)}=\prod_{x\in X} d_x!,
\]
以及规范体积密度
\[
g(f):=2^{-n}\log\card{\mathcal G(f)}=2^{-n}\sum_{x\in X}\log(d_x!).
\]
再定义界面面积（切割边数）
\[
A(f):=\card{\{\{u,v\}\in E_n:\ f(u)\neq f(v)\}}.
\]
\end{definition}

\begin{theorem}[Shannon--Stirling 双展开与体--界熵缺口恒等式]\label{thm:hypercube-gauge-volume-shannon-stirling-gap}
\leanpartial{paper\_hypercube\_gauge\_volume\_shannon\_stirling\_gap}{the requested Lean signature is false for the current `Omega.Folding.binGaugeG` normalization: with \(m=1\), \(|\alpha|=2\), and \(d(a)\equiv 1\), the left side is \(0\) while the displayed right side is already about \(1.3052\) before adding the constrained remainder \(R\ge 0\)}
在定义 \ref{def:hypercube-coarsegraining-area-gauge-volume} 的设定下，令推前分布 \(p_x:=d_x/2^n\)（\(x\in X\)），并记 Shannon 熵
\[
H(p):=-\sum_{x\in X}p_x\log p_x.
\]
则存在显式余项 \(R(f)\ge 0\) 使得严格恒等式成立：
\[
\boxed{\;
g(f)=n\log 2-1-H(p)
\;+\;\frac1{2^{n+1}}\sum_{x\in X}\log(2\pi d_x)
\;+\;R(f).
\;}
\]
并且余项满足统一上界
\[
0\le R(f)\le \frac{|X|}{12\cdot 2^n}.
\]
因此
\[
\bigl|g(f)-(n\log 2-1-H(p))\bigr|
\le
\frac{|X|}{2^{n+1}}\log(2\pi 2^n)+\frac{|X|}{12\cdot 2^n}.
\]
特别地，若 \(|X|/2^n\to 0\)（当 \(n\to\infty\)），则
\(
g(f)=n\log 2-1-H(p)+o(1)
\)。
\end{theorem}
\begin{proof}
对整数 \(m\ge 1\) 取带余项的 Stirling 展开
\[
\log(m!)=m\log m-m+\tfrac12\log(2\pi m)+\theta_m,
\qquad
0\le \theta_m\le \frac1{12m}.
\]
将其代入 \(g(f)=2^{-n}\sum_x\log(d_x!)\) 并逐项求和得
\[
g(f)
=2^{-n}\sum_x\bigl(d_x\log d_x-d_x\bigr)
\;+\;\frac1{2^{n+1}}\sum_x\log(2\pi d_x)
\;+\;2^{-n}\sum_x\theta_{d_x}.
\]
注意 \(\sum_x d_x=2^n\)，故 \(2^{-n}\sum_x d_x=1\)，并且
\[
2^{-n}\sum_x d_x\log d_x
=\sum_x p_x\log(2^n p_x)
=n\log 2+\sum_x p_x\log p_x
=n\log 2-H(p).
\]
令
\(
R(f):=2^{-n}\sum_x\theta_{d_x}
\)，则 \(R(f)\ge 0\) 且
\[
0\le R(f)\le 2^{-n}\sum_x\frac1{12d_x}\le \frac{|X|}{12\cdot 2^n}.
\]
从而得到所述恒等式与误差界。证毕。
\end{proof}

\begin{proposition}[规范体积密度的 Rényi 型上界（由碰撞矩直接控制）]\label{prop:hypercube-gauge-volume-renyi-upperbound}
\leanverified{paper\_hypercube\_gauge\_volume\_renyi\_upperbound}
在定义 \ref{def:hypercube-coarsegraining-area-gauge-volume} 的设定下，令推前分布 \(p_x:=d_x/2^n\)（\(x\in X\)）。
对任意 \(q>1\) 定义 \(q\)-碰撞矩
\[
S_q(f):=\sum_{x\in X} d_x^q.
\]
则规范体积密度满足上界
\begin{equation}\label{eq:hypercube-gauge-volume-renyi-upperbound}
\boxed{\;
g(f)\ \le\ \frac{1}{q-1}\Bigl(\log S_q(f)-n\log 2\Bigr)\ -\ 1\ +\ \frac{|X|}{2^n}.\;}
\end{equation}
\end{proposition}
\begin{proof}
由基本不等式 \(\log(m!)\le m\log m-m+1\)（\(m\ge 1\)）得
\[
g(f)=2^{-n}\sum_{x\in X}\log(d_x!)
\le 2^{-n}\sum_{x\in X}\bigl(d_x\log d_x-d_x+1\bigr)
=\sum_{x\in X}p_x\log d_x-1+\frac{|X|}{2^n}.
\]
又对 \(q>1\)，由 Jensen 不等式（\(\log\) 凹）得到
\[
\sum_{x\in X}p_x\log d_x
=\frac{1}{q-1}\sum_{x\in X}p_x\log d_x^{\,q-1}
\le \frac{1}{q-1}\log\sum_{x\in X}p_x d_x^{\,q-1}
=\frac{1}{q-1}\log\Bigl(\frac{S_q(f)}{2^n}\Bigr),
\]
代回即得 \eqref{eq:hypercube-gauge-volume-renyi-upperbound}。证毕。
\end{proof}
\leanpartial{paper\_hypercube\_gauge\_volume\_renyi\_upperbound}{the requested Lean signature is false: with \(X=\{\ast\}\), \(d(\ast)=2\), \(n=1\), \(q=2\), the left side is \(\log 2 / 2\) while the right side is \(\log 2 - 1/2\)}

\begin{proposition}[离散全息斯托克不等式：规范体积密度的切割下界]\label{prop:hypercube-discrete-holographic-stokes-gauge-area}
在定义 \ref{def:hypercube-coarsegraining-area-gauge-volume} 的设定下，恒有
\[
\boxed{\;
g(f)\ \ge\ (n\log 2-1)\ -\ 2\log 2\cdot\frac{A(f)}{2^n}\ +\ \frac{\card{X}}{2^n}.
\;}
\]
\leanverified{paper\_hypercube\_discrete\_holographic\_stokes\_gauge\_area}
\end{proposition}

\begin{proof}
\textbf{（阶乘下界）}
对整数 \(m\ge 1\) 有
\[
\log(m!)=\sum_{k=1}^{m}\log k\ \ge\ \int_{1}^{m}\log t\,\dd t\ =\ m\log m-m+1.
\]
故
\[
\sum_{x\in X}\log(d_x!)
\ge
\sum_{x\in X}(d_x\log d_x-d_x+1)
=
\sum_{x\in X}d_x\log d_x-2^n+\card{X}.
\]

\textbf{（超立方体边缘等周不等式）}
对任意 \(S\subseteq \Omega_n\)，记其边界边集
\(
\partial S:=\{\{u,v\}\in E_n:\ u\in S,\ v\notin S\}
\)。
超立方体的标准边缘等周不等式给出
\[
\card{\partial S}
\ge
\card{S}\log_2\!\Bigl(\frac{2^n}{\card{S}}\Bigr)
=
\card{S}\bigl(n-\log_2\card{S}\bigr).
\]
应用于每个纤维 \(S_x\) 得
\[
\card{\partial S_x}\ge d_x\bigl(n-\log_2 d_x\bigr)
\quad\Longrightarrow\quad
d_x\log_2 d_x\ \ge\ nd_x-\card{\partial S_x}.
\]
对 \(x\in X\) 求和并用 \(\sum_x d_x=2^n\)，得
\[
\sum_{x\in X}d_x\log_2 d_x
\ge
n2^n-\sum_{x\in X}\card{\partial S_x}.
\]
注意每条切割边 \(\{u,v\}\)（满足 \(f(u)\neq f(v)\)）恰被计入 \(\sum_x\card{\partial S_x}\) 两次，故 \(\sum_x\card{\partial S_x}=2A(f)\)。
于是
\[
\sum_{x\in X}d_x\log d_x
=
\log 2\cdot \sum_{x\in X}d_x\log_2 d_x
\ge
\log 2\cdot (n2^n-2A(f)).
\]

\textbf{（合并）}
代回阶乘下界并除以 \(2^n\) 得
\[
g(f)
\ge
2^{-n}\Bigl(\log 2\cdot (n2^n-2A(f))-2^n+\card{X}\Bigr),
\]
即所示不等式。证毕。
\end{proof}

\begin{theorem}[全息 Kraft--Gödel 侧信息长度下界]\label{thm:hypercube-kraft-godel-sideinfo-lower-bound}
\leanverified{paper\_hypercube\_kraft\_godel\_sideinfo\_lower\_bound}
令 \(f:\Omega_n=\{0,1\}^n\twoheadrightarrow X\) 为满射粗粒化，并沿用界面面积 \(A(f)\) 与规范体积密度 \(g(f)\) 的记号（定义 \ref{def:hypercube-coarsegraining-area-gauge-volume}）。
固定字母表 \(\Sigma\)，\(\card{\Sigma}=M\ge 2\)。
设存在映射
\[
\Phi:\Omega_n\longrightarrow X\times \Sigma^\ast
\]
满足 \(\Phi(\omega)=\bigl(f(\omega),\Phi_2(\omega)\bigr)\) 且 \(\Phi\) 在 \(\Omega_n\) 上单射，并且对每个 \(x\in X\)，集合 \(\{\Phi_2(\omega):\omega\in f^{-1}(x)\}\subset \Sigma^\ast\) 为前缀码。
定义平均侧信息长度
\[
L(\Phi):=2^{-n}\sum_{\omega\in\Omega_n}\bigl|\Phi_2(\omega)\bigr|.
\]
则有下界
\[
\boxed{\;
L(\Phi)\ \ge\ \frac{g(f)+1-\card{X}/2^n}{\log M}
\ \ge\ \frac{n\log 2-\frac{\log 2}{2^{n-1}}A(f)}{\log M}.\;}
\]
特别地，当 \(M=2\) 时，
\[
L(\Phi)\ \ge\ n-\frac{A(f)}{2^{n-1}}.
\]
\end{theorem}
\begin{proof}
对固定 \(x\in X\) 记纤维大小 \(d_x:=\card{f^{-1}(x)}\)，并将纤维上的码字长度记为 \(\ell_{x,1},\dots,\ell_{x,d_x}\)。
由于 \(\{\Phi_2(\omega):\omega\in f^{-1}(x)\}\) 为 \(M\)-元前缀码，Kraft 不等式给出
\[
\sum_{j=1}^{d_x}M^{-\ell_{x,j}}\le 1.
\]
注意 \(t\mapsto M^{-t}=\exp(-t\log M)\) 为凸函数，Jensen 不等式得到
\[
\frac1{d_x}\sum_{j=1}^{d_x}M^{-\ell_{x,j}}
\ \ge\
M^{-\frac1{d_x}\sum_{j=1}^{d_x}\ell_{x,j}}
=M^{-L_x},
\qquad
L_x:=\frac1{d_x}\sum_{j=1}^{d_x}\ell_{x,j}.
\]
从而
\[
1\ \ge\ \sum_{j=1}^{d_x}M^{-\ell_{x,j}}\ \ge\ d_x\,M^{-L_x},
\qquad\text{即}\qquad
L_x\ge \log_M d_x=\frac{\log d_x}{\log M}.
\]
令 \(p_x:=d_x/2^n\)，并对 \(x\) 加权求和得
\[
L(\Phi)=\sum_{x\in X}p_x L_x\ \ge\ \frac{1}{\log M}\sum_{x\in X}p_x\log d_x.
\]
另一方面，对整数 \(m\ge 1\) 有基本上界 \(\log(m!)\le m\log m-m+1\)，故
\[
g(f)=2^{-n}\sum_{x\in X}\log(d_x!)
\le
2^{-n}\sum_{x\in X}\bigl(d_x\log d_x-d_x+1\bigr)
=\sum_{x\in X}p_x\log d_x-1+\frac{\card{X}}{2^n},
\]
即 \(\sum_x p_x\log d_x\ge g(f)+1-\card{X}/2^n\)。
合并得到第一条不等式。
最后由命题 \ref{prop:hypercube-discrete-holographic-stokes-gauge-area}，
\[
g(f)\ \ge\ (n\log 2-1)-2\log 2\cdot\frac{A(f)}{2^n}+\frac{\card{X}}{2^n},
\]
从而
\[
g(f)+1-\frac{\card{X}}{2^n}\ \ge\ n\log 2-\frac{\log 2}{2^{n-1}}A(f),
\]
代回得到第二条不等式。证毕。
\end{proof}

\endinput


\begin{proposition}[面积律的熵上界与二相刚性]\label{prop:hypercube-entropy-area-rigidity}
\leanverified{paper\_hypercube\_entropy\_area\_rigidity}
在定义 \ref{def:hypercube-coarsegraining-area-gauge-volume} 的设定下，令 \(p_x:=d_x/2^n\) 为纤维分布，并记 Shannon 熵
\[
H(p):=-\sum_{x\in X}p_x\log p_x.
\]
则有：
\begin{enumerate}
  \item \label{it:hypercube-entropy-area} 熵--面积上界
  \[
  \boxed{\;
  H(p)\ \le\ 2\log 2\cdot \frac{A(f)}{2^n}.
  \;}
  \]
  \item \label{it:hypercube-area-collapse} 若 \(A(f)=o(2^n)\)（当 \(n\to\infty\)），则 \(\max_x p_x\to 1\)。
  \item \label{it:hypercube-area-saturation} 记 \(\card{E_n}=n2^{n-1}\)。若 \(A(f)=\card{E_n}-o(2^n)\)，则对均匀随机边 \(\{U,V\}\in E_n\) 有 \(\PP(f(U)=f(V))\to 0\)。
\end{enumerate}
\end{proposition}

\begin{proof}
\textbf{（\ref{it:hypercube-entropy-area}）}
由边缘等周不等式，\(\card{\partial S_x}\ge d_x\log_2(2^n/d_x)=d_x\log_2(1/p_x)\)。
求和得
\[
\sum_{x\in X}\card{\partial S_x}
\ge
2^n\sum_{x\in X}p_x\log_2(1/p_x)
=
2^n H_2(p),
\]
其中 \(H_2(p):=-\sum_x p_x\log_2 p_x\) 为以 \(2\) 为底的熵。
又 \(\sum_x\card{\partial S_x}=2A(f)\)，故 \(H(p)=\log 2\cdot H_2(p)\le 2\log 2\cdot A(f)/2^n\)。

\textbf{（\ref{it:hypercube-area-collapse}）}
设 \(p_{\max}=1-\delta\)。则对任意分布有下界 \(H(p)\ge h(\delta)\)，其中
\(
h(\delta):=-\delta\log\delta-(1-\delta)\log(1-\delta)
\)。
由 \ref{it:hypercube-entropy-area} 若 \(A(f)=o(2^n)\) 则 \(H(p)\to 0\)，从而 \(\delta\to 0\)，即 \(p_{\max}\to 1\)。

\textbf{（\ref{it:hypercube-area-saturation}）}
记内部边数
\[
I(f):=\card{\{\{u,v\}\in E_n:\ f(u)=f(v)\}},
\]
则 \(A(f)=\card{E_n}-I(f)\)。
若 \(A(f)=\card{E_n}-o(2^n)\)，则 \(I(f)=o(2^n)\)。
对均匀随机边 \(\{U,V\}\) 有
\[
\PP(f(U)=f(V))=\frac{I(f)}{\card{E_n}}=\frac{o(2^n)}{n2^{n-1}}\to 0.
\]
证毕。
\end{proof}

\begin{definition}[边界邻接多图]\label{def:hypercube-boundary-multigraph}
在定义 \ref{def:hypercube-coarsegraining-area-gauge-volume} 的设定下，定义边界邻接多图 \(\mathsf B_f\) 为：其顶点集为 \(X\)，并对每条切割边 \(\{u,v\}\in E_n\)（满足 \(f(u)\neq f(v)\)）添加一条连接 \(f(u)\) 与 \(f(v)\) 的无向边（允许重边与自环；本构造中自环不出现）。
于是 \(\card{E(\mathsf B_f)}=A(f)\)。
\end{definition}

\begin{proposition}[回路圆维的欧拉恒等式]\label{prop:hypercube-boundary-multigraph-h1-cdim}
\leanverified{paper\_hypercube\_boundary\_multigraph\_h1\_cdim}
在定义 \ref{def:hypercube-boundary-multigraph} 的设定下，记 \(\mathsf B_f\) 的连通分量数为 \(c(\mathsf B_f)\)。
则其一维同调为自由阿贝尔群，且
\[
H_1(\mathsf B_f;\ZZ)\ \cong\ \ZZ^{\,A(f)-\card{X}+c(\mathsf B_f)}.
\]
因此（按圆维在有限生成交换群上的自由秩口径）
\[
\boxed{\;
\cdim\bigl(H_1(\mathsf B_f;\ZZ)\bigr)=A(f)-\card{X}+c(\mathsf B_f).
\;}
\]
\end{proposition}

\begin{proof}
将 \(\mathsf B_f\) 视作一维 CW 复形。其链复形满足 \(C_1\cong \ZZ^{\card{E}}\)、\(C_0\cong \ZZ^{\card{V}}\)，且
\(
\operatorname{rank}H_1=\card{E}-\card{V}+c
\)。
代入 \(\card{E}=A(f)\)、\(\card{V}=\card{X}\) 即得秩公式，并得同构 \(H_1(\mathsf B_f;\ZZ)\cong \ZZ^{\operatorname{rank}H_1}\)。
最后由 \(\cdim(\ZZ^r)=r\) 得圆维等式。证毕。
\end{proof}

% Coarse-graining "c-theorem" for cycle rank under contracting connected blocks.

\begin{theorem}[局部粗粒化的回路圆维分解：严格单调性与损失恒等式]\label{thm:graph-coarsegraining-cycle-rank-decomposition}
\leanverified{paper\_graph\_coarsegraining\_cycle\_rank\_decomposition}
令 \(G\) 为有限无自环多重图（允许重边），其顶点集为 \(V(G)\)、边集为 \(E(G)\)，连通分量数记为 \(c(G)\)。
设 \(\pi:V(G)\twoheadrightarrow Y\) 为满射，并满足：
\begin{enumerate}
  \item 对任意 \(y\in Y\)，纤维 \(\pi^{-1}(y)\) 包含于 \(G\) 的某个连通分量内；
  \item 对任意 \(y\in Y\)，诱导子图 \(G[\pi^{-1}(y)]\) 连通。
\end{enumerate}
定义商图 \(G/\pi\) 为将每个纤维 \(\pi^{-1}(y)\) 收缩为单点并删除纤维内部边后所得之多重图（仅保留连接不同纤维的边，并继承其多重性）。
则 \(H_1(G;\ZZ)\) 与 \(H_1(G/\pi;\ZZ)\) 皆为自由阿贝尔群，且回路圆维满足严格恒等式
\[
\cdim\,H_1(G/\pi;\ZZ)
=\cdim\,H_1(G;\ZZ)\ -\ \sum_{y\in Y}\cdim\,H_1\!\bigl(G[\pi^{-1}(y)];\ZZ\bigr).
\]
因此
\[
\cdim\,H_1(G/\pi;\ZZ)\le \cdim\,H_1(G;\ZZ),
\]
并且等号成立当且仅当对每个 \(y\in Y\)，诱导子图 \(G[\pi^{-1}(y)]\) 为树（亦即其回路圆维为 \(0\)）。
\end{theorem}
\begin{proof}
设 \(|V(G)|=V\)、\(|E(G)|=E\)，并记 \(|Y|=K\)。
令 \(E_{\mathrm{in}}\) 为纤维内部边数，\(E_{\mathrm{out}}:=E-E_{\mathrm{in}}\) 为跨纤维边数，则由定义 \(G/\pi\) 的边数为 \(E_{\mathrm{out}}\)，顶点数为 \(K\)。
由条件 (1) 可知收缩不跨连通分量，故 \(c(G/\pi)=c(G)\)。
对任意有限图，一维同调的秩满足欧拉恒等式
\[
\cdim\,H_1(G;\ZZ)=E-V+c(G),
\qquad
\cdim\,H_1(G/\pi;\ZZ)=E_{\mathrm{out}}-K+c(G).
\]
因此
\[
\cdim\,H_1(G/\pi;\ZZ)
=\cdim\,H_1(G;\ZZ)+(V-K)-E_{\mathrm{in}}.
\]
另一方面，对每个 \(y\in Y\)，由于 \(G[\pi^{-1}(y)]\) 连通，记其顶点数为 \(V_y:=|\pi^{-1}(y)|\)、边数为 \(E_y\)，则
\[
\cdim\,H_1\!\bigl(G[\pi^{-1}(y)];\ZZ\bigr)=E_y-V_y+1,
\qquad\text{即}\qquad
E_y=V_y-1+\cdim\,H_1\!\bigl(G[\pi^{-1}(y)];\ZZ\bigr).
\]
对 \(y\) 求和并用 \(\sum_y V_y=V\)、\(\sum_y E_y=E_{\mathrm{in}}\) 得
\[
E_{\mathrm{in}}=(V-K)+\sum_{y\in Y}\cdim\,H_1\!\bigl(G[\pi^{-1}(y)];\ZZ\bigr).
\]
代回前式即得所示恒等式。
末句由每一项非负且为零当且仅当相应诱导子图为树得到。证毕。
\end{proof}

\endinput


\begin{theorem}[加权回路格判别式的矩阵树恒等式]\label{thm:graph-cycle-lattice-weighted-discriminant}
\leanverified{paper\_graph\_cycle\_lattice\_weighted\_discriminant}
令 \(G=(V,E)\) 为有限连通无自环多重图，固定一组取向并把 \(1\)-链识别为 \(\ZZ^{E}\)。
令 \(\partial:\ZZ^{E}\to \ZZ^{V}\) 为边界算子，并令
\[
\Lambda(G):=\ker(\partial)\subset \ZZ^{E},
\qquad
r:=\rank_{\ZZ}\Lambda(G)=\rank_{\ZZ}H_1(G;\ZZ)=|E|-|V|+1.
\]
给定正权 \(\rho:E\to(0,\infty)\)，定义二次型
\[
Q_\rho(z):=\sum_{e\in E}\rho_e\,z_e^2,\qquad z=(z_e)_{e\in E}\in\RR^{E}.
\]
取任意 \(\ZZ\)-基 \(c_1,\dots,c_r\) 生成 \(\Lambda(G)\)，并令 \(C\in\ZZ^{r\times |E|}\) 为基矩阵（第 \(j\) 行为 \(c_j^{\mathsf T}\)）。
定义
\[
\Gamma_\rho:=C\,\mathrm{diag}(\rho_e)\,C^{\mathsf T}\in \mathrm{Sym}_r^+(\RR).
\]
则 \(\det(\Gamma_\rho)\) 与基的选择无关，且满足精确恒等式
\[
\boxed{\;
\det(\Gamma_\rho)=\Bigl(\prod_{e\in E}\rho_e\Bigr)\,\tau_{1/\rho}(G),
\qquad
\tau_{1/\rho}(G):=\sum_{T\ \mathrm{spanning\ tree}}\ \prod_{e\in T}\rho_e^{-1}.
\;}
\]
\end{theorem}
\begin{proof}
基不变性来自：更换 \(\ZZ\)-基等价于 \(C\mapsto UC\)（\(U\in\GL_r(\ZZ)\)），从而 \(\Gamma_\rho\mapsto U\Gamma_\rho U^{\mathsf T}\)，其行列式不变。
\par
取 \(G\) 的一棵生成树 \(T_0\subseteq E\) 并由此构造基本回路基，使得存在 unimodular 矩阵 \(M\in\GL_{|E|}(\ZZ)\)，其上方 \(|V|-1\) 行为删去一行后的约化关联矩阵 \(\tilde B\in\ZZ^{(|V|-1)\times |E|}\)，下方 \(r\) 行恰为某个回路基矩阵 \(C\)，从而 \(\tilde B C^{\mathsf T}=0\)。
\par
对 \(\Gamma_\rho=C\,\mathrm{diag}(\rho_e)\,C^{\mathsf T}\) 应用 Cauchy--Binet 公式，得
\[
\det(\Gamma_\rho)
=\sum_{S\subseteq E,\ |S|=r}\det(C_S)^2\prod_{e\in S}\rho_e,
\]
其中 \(C_S\) 为取列集 \(S\) 的 \(r\times r\) 子矩阵。
由于 \(M\) unimodular，其互补子式满足
\(
\det(C_S)=\pm\det(\tilde B_{E\setminus S})
\)，
其中 \(\tilde B_{E\setminus S}\) 为 \(\tilde B\) 取列集 \(E\setminus S\) 的 \((|V|-1)\times(|V|-1)\) 子矩阵。
由矩阵树定理的标准整数形式（关联矩阵全酉模性质），\(\det(\tilde B_{E\setminus S})=\pm 1\) 当且仅当 \(E\setminus S\) 构成生成树，否则为 \(0\)。
因此
\[
\det(\Gamma_\rho)
=\sum_{T\ \mathrm{spanning\ tree}}\ \prod_{e\in E\setminus T}\rho_e
=\Bigl(\prod_{e\in E}\rho_e\Bigr)\sum_{T}\prod_{e\in T}\rho_e^{-1}
=\Bigl(\prod_{e\in E}\rho_e\Bigr)\tau_{1/\rho}(G).
\]
证毕。
\end{proof}

\begin{proposition}[生成树--回路判别式的相位化：圆维、体积常数与 Lee--Yang 圆谱]\label{prop:graph-cycle-lattice-leyang-phase-polynomial}
令 \(G=(V,E)\) 为有限连通无自环多重图，并取权重 \(\rho:E\to(0,\infty)\)。
令 \(L_{1/\rho}\) 为以导通率 \(c_e:=\rho_e^{-1}\) 加权的 Laplacian，并记其谱为
\[
0=\lambda_1<\lambda_2\le\cdots\le \lambda_{|V|}.
\]
取任意 \(D\ge \lambda_{|V|}/2\)，并对 \(j\ge 2\) 定义相位角 \(\theta_j\in(0,\pi]\) 使得
\[
\cos\theta_j=1-\frac{\lambda_j}{2D}.
\]
定义 Lee--Yang 相位多项式
\[
Q_{G,D}(z):=\prod_{j=2}^{|V|}(z-e^{\mathrm{i}\theta_j})(z-e^{-\mathrm{i}\theta_j})
=\prod_{j=2}^{|V|}\bigl(z^2-2\cos\theta_j\,z+1\bigr).
\]
则 \(Q_{G,D}\) 的全部非平凡零点均落在单位圆上，且加权生成树多项式值满足显式相位分解
\[
\boxed{\;
\tau_{1/\rho}(G)
=\frac{1}{|V|}\prod_{j=2}^{|V|}\lambda_j
=\frac{(4D)^{|V|-1}}{|V|}\prod_{j=2}^{|V|}\sin^2\!\Bigl(\frac{\theta_j}{2}\Bigr).\;}
\]
从而与定理 \ref{thm:graph-cycle-lattice-weighted-discriminant} 合并得到
\[
\det(\Gamma_\rho)
=\Bigl(\prod_{e\in E}\rho_e\Bigr)\tau_{1/\rho}(G)
=\Bigl(\prod_{e\in E}\rho_e\Bigr)\frac{(4D)^{|V|-1}}{|V|}\prod_{j=2}^{|V|}\sin^2\!\Bigl(\frac{\theta_j}{2}\Bigr).
\]
\end{proposition}
\begin{proof}
由加权 Kirchhoff 矩阵树定理，
\(
\tau_{1/\rho}(G)=\det\bigl((L_{1/\rho})_{\mathrm{red}}\bigr)=\frac1{|V|}\prod_{j=2}^{|V|}\lambda_j
\)。
由 \(\cos\theta_j=1-\lambda_j/(2D)\) 得
\[
\lambda_j=2D(1-\cos\theta_j)=4D\sin^2(\theta_j/2),
\]
对 \(j=2,\dots,|V|\) 相乘即得相位分解式。
多项式 \(Q_{G,D}\) 的零点定义为 \(\{e^{\pm i\theta_j}\}_{j=2}^{|V|}\)，故其非平凡零点全落在单位圆上。
最后一式由定理 \ref{thm:graph-cycle-lattice-weighted-discriminant} 的 \(\det(\Gamma_\rho)=(\prod_e\rho_e)\tau_{1/\rho}(G)\) 直接得到。证毕。
\end{proof}

\endinput



% Modular inversion of weighted theta series on integer cycle lattices.

\begin{theorem}[回路格加权 \texorpdfstring{$\Theta$}{Theta} 函数的模反演公式与一维椭圆曲线接口]\label{thm:graph-cycle-lattice-theta-modular-inversion}
\leanverified{paper\_graph\_cycle\_lattice\_theta\_modular\_inversion}
令 \(G=(V,E)\) 为有限连通无自环多重图，固定一组取向并把 \(1\)-链识别为 \(\ZZ^{E}\)。
令 \(\partial:\ZZ^{E}\to \ZZ^{V}\) 为边界算子，并令回路格
\[
\Lambda(G):=\ker(\partial)\subset \ZZ^{E},
\qquad
r:=\rank_{\ZZ}\Lambda(G)=\rank_{\ZZ}H_1(G;\ZZ).
\]
给定正权 \(\rho:E\to(0,\infty)\)，定义二次型
\[
Q_\rho(z):=\sum_{e\in E}\rho_e\,z_e^2,
\qquad z=(z_e)_{e\in E}.
\]
定义加权 \(\Theta\) 函数
\[
\Theta_{G,\rho}(t):=\sum_{\psi\in\Lambda(G)}\exp\!\bigl(-\pi t\,Q_\rho(\psi)\bigr),
\qquad t>0.
\]
取任意 \(\ZZ\)-基 \(c_1,\dots,c_r\) 生成 \(\Lambda(G)\)，并令 \(C\in\ZZ^{r\times |E|}\) 为基矩阵（第 \(j\) 行为 \(c_j^{\mathsf T}\)）。
记
\[
\Gamma_\rho:=C\,\mathrm{diag}(\rho_e)\,C^{\mathsf T}\in\mathrm{Sym}_r^+(\RR).
\]
则存在对偶格 \(\Lambda(G)^\vee\subset \Lambda(G)_{\RR}^\ast\) 使得对一切 \(t>0\) 成立精确恒等式
\[
\boxed{\;
\Theta_{G,\rho}(t)
=t^{-r/2}\,(\det\Gamma_\rho)^{-1/2}
\sum_{\xi\in \Lambda(G)^\vee}\exp\!\left(-\pi\,\frac{\langle\xi,\Gamma_\rho^{-1}\xi\rangle}{t}\right).
\;}
\]
从而当 \(t\downarrow 0\) 时有主项渐近
\[
\Theta_{G,\rho}(t)\sim t^{-r/2}(\det\Gamma_\rho)^{-1/2},
\qquad
\Theta_{G,\rho}(t)\to 1\ (t\to\infty).
\]
当 \(r=1\) 时，\(\Lambda(G)=\ZZ\cdot v\) 且 \(Q_\rho(v)=:L>0\)，于是
\[
\Theta_{G,\rho}(t)=\sum_{k\in\ZZ}e^{-\pi t L k^2},
\]
其模反演对应于 Jacobi \(\theta\) 函数在参数 \(\tau=i t L\) 下的 \(\tau\mapsto-1/\tau\) 变换，并自然给出椭圆曲线
\(
E_\tau=\CC/(\ZZ+\tau\ZZ)
\)
作为能量尺度 \(tL\) 的周期模数接口。
\end{theorem}
\begin{proof}
任意 \(\psi\in\Lambda(G)\) 可唯一写为 \(\psi=\sum_{j=1}^r k_j c_j\)（\(k\in\ZZ^r\)），从而
\[
Q_\rho(\psi)
=(C^{\mathsf T}k)^{\mathsf T}\mathrm{diag}(\rho_e)(C^{\mathsf T}k)
=k^{\mathsf T}\Gamma_\rho k.
\]
因此
\(
\Theta_{G,\rho}(t)=\sum_{k\in\ZZ^r}e^{-\pi t\,k^{\mathsf T}\Gamma_\rho k}
\)。
对 Schwartz 高斯函数 \(f_t(x):=\exp(-\pi t\,x^{\mathsf T}\Gamma_\rho x)\) 在 \(\RR^r\) 上应用 Poisson 求和，
\[
\sum_{k\in\ZZ^r}f_t(k)=\sum_{m\in\ZZ^r}\widehat f_t(m),
\]
并用高斯 Fourier 变换
\(
\widehat f_t(m)=t^{-r/2}(\det\Gamma_\rho)^{-1/2}\exp(-\pi\,m^{\mathsf T}\Gamma_\rho^{-1}m/t)
\)
即可得到所述恒等式；将右端的 \(\ZZ^r\) 识别为 \(\Lambda(G)^\vee\) 在对偶基下的坐标表达即得基无关表述。
极限行为由右端 \(m=0\) 主项与支配收敛立得。
当 \(r=1\) 时取 \(L:=\Gamma_\rho\) 即得标准 Jacobi 形式及其模变换接口。
证毕。
\end{proof}

\begin{corollary}[仿射回路扭量的 \texorpdfstring{$\Theta$}{Theta} 函数与相位对偶]\label{cor:graph-cycle-affine-torsor-theta-poisson-phase}
\leanverified{paper\_graph\_cycle\_affine\_torsor\_theta\_poisson\_phase}
沿用定理 \ref{thm:graph-cycle-lattice-theta-modular-inversion} 的记号，并令 \(\Lambda:=\Lambda(G)\subset\Lambda_{\RR}:=\Lambda(G)\otimes_{\ZZ}\RR\)。
对任意平移参数 \(\psi_0\in\Lambda_{\RR}\)，定义仿射 \(\Theta\) 函数
\[
\Theta_{G,\rho}^{(\psi_0)}(t):=\sum_{\psi\in\Lambda}\exp\!\bigl(-\pi t\,Q_\rho(\psi+\psi_0)\bigr),
\qquad t>0.
\]
则存在对偶格 \(\Lambda^\vee\subset \Lambda_{\RR}^\ast\)（同定理 \ref{thm:graph-cycle-lattice-theta-modular-inversion}）使得对一切 \(t>0\) 成立精确恒等式
\[
\boxed{\;
\Theta_{G,\rho}^{(\psi_0)}(t)
=t^{-r/2}\,(\det\Gamma_\rho)^{-1/2}
\sum_{\xi\in \Lambda^\vee}\exp\!\left(-\pi\,\frac{\langle\xi,\Gamma_\rho^{-1}\xi\rangle}{t}\right)\,
e^{2\pi i\langle\xi,\psi_0\rangle}.
\;}
\]
其中相位因子 \(e^{2\pi i\langle\xi,\psi_0\rangle}\) 仅由仿射扭量类 \(\psi_0\ \mathrm{mod}\ \Lambda\) 决定。
\end{corollary}
\begin{proof}
令 \(f_t(x):=\exp(-\pi t\,x^{\mathsf T}\Gamma_\rho x)\)（沿定理 \ref{thm:graph-cycle-lattice-theta-modular-inversion} 的坐标化）。
则
\(
\Theta_{G,\rho}^{(\psi_0)}(t)=\sum_{k\in\ZZ^r} f_t(k+k_0)
\)，其中 \(k_0\in\RR^r\) 为 \(\psi_0\) 在所选基下的坐标。
对平移后的 Poisson 求和公式
\[
\sum_{k\in\ZZ^r} f_t(k+k_0)=\sum_{m\in\ZZ^r}\widehat f_t(m)\,e^{2\pi i\langle m,k_0\rangle}
\]
应用高斯 Fourier 变换 \(\widehat f_t(m)=t^{-r/2}(\det\Gamma_\rho)^{-1/2}\exp(-\pi m^{\mathsf T}\Gamma_\rho^{-1}m/t)\)，并将 \(m\) 识别为对偶格元 \(\xi\in\Lambda^\vee\) 即得所示恒等式。
由于 \(k_0\) 改变一个整数向量仅使相位乘以 \(1\)，故相位只依赖于 \(\psi_0\ \mathrm{mod}\ \Lambda\)。证毕。
\end{proof}

\endinput


\begin{theorem}[边界回路格的 \texorpdfstring{$\Theta$}{Theta} 函数极限律：圆维主项与树数校正]\label{thm:hypercube-cycle-lattice-theta-asymptotic}
\leanverified{paper\_hypercube\_cycle\_lattice\_theta\_asymptotic}
令 \(G=(V,E)\) 为有限连通多重图，固定一组取向并将 \(1\)-链视作 \(\ZZ^{E}\subset\RR^{E}\)。
令 \(B\in\ZZ^{V\times E}\) 为取向关联矩阵，并令 \(\tilde B\in\ZZ^{(|V|-1)\times E}\) 为删去一行后的约化关联矩阵。
定义回路格
\[
\Lambda(G):=\ker(\tilde B)\cap\ZZ^{E},
\qquad
r:=\operatorname{rank}_{\ZZ}\Lambda(G)=|E|-|V|+1.
\]
以 \(\RR^{E}\) 的标准内积诱导 Euclid 范数 \(|\cdot|\)，并定义回路格的 theta 函数
\[
\Theta_G(t):=\sum_{z\in\Lambda(G)}e^{-\pi t|z|^2},
\qquad t>0.
\]
记 \(\tau(G)\) 为 \(G\) 的生成树个数（Kirchhoff 复杂度）。
则存在常数 \(c(G)>0\) 使得当 \(t\downarrow 0\) 时有渐近展开
\[
\boxed{\;
\Theta_G(t)=t^{-r/2}\,\tau(G)^{-1/2}\Bigl(1+O\bigl(e^{-c(G)/t}\bigr)\Bigr).
\;}
\]
从而
\[
\boxed{\;
\log\Theta_G(t)=\frac r2\log\frac1t-\frac12\log\tau(G)+O\bigl(e^{-c(G)/t}\bigr).
\;}
\]
特别地，取 \(G=\mathsf B_f\) 时，\(r=\cdim(H_1(\mathsf B_f;\ZZ))\)（命题 \ref{prop:hypercube-boundary-multigraph-h1-cdim}），故上述自由能主项系数即为边界回路圆维。
\end{theorem}
\begin{proof}
\textbf{（i）回路格协体积与树数）}
选取 \(G\) 的一棵生成树 \(T\subseteq E\)。
记 \(E= T\sqcup C\) 为边集分解，其中 \(|T|=|V|-1\)、\(|C|=r\)。
按该分解将 \(\tilde B\) 分块写成
\(
\tilde B=[\tilde B_T\ \ \tilde B_C]
\)。
由于 \(T\) 为树，\(\tilde B_T\in\ZZ^{(|V|-1)\times(|V|-1)}\) 可逆且 \(\det(\tilde B_T)=\pm 1\)（关联矩阵的树子式为 \(\pm 1\)）。
于是对任意 \(u\in\ZZ^{r}\)，向量
\[
z(u):=\begin{pmatrix}
 -\tilde B_T^{-1}\tilde B_C\,u\\
 u
\end{pmatrix}\in\ZZ^{E}
\]
满足 \(\tilde B z(u)=0\)，并且 \(u\mapsto z(u)\) 给出 \(\Lambda(G)\) 的一组 \(\ZZ\)-基。
令 \(Z\in\ZZ^{E\times r}\) 为该基矩阵，则回路格的协体积平方为
\[
\operatorname{covol}(\Lambda(G))^2=\det(Z^\top Z).
\]
写 \(X:=-\tilde B_T^{-1}\tilde B_C\)，则 \(Z=\binom{X}{I_r}\)，从而
\[
\det(Z^\top Z)=\det(I_r+X^\top X)=\det(I_{|V|-1}+XX^\top),
\]
其中最后一步用 Sylvester 行列式恒等式。
进一步
\[
XX^\top=\tilde B_T^{-1}\tilde B_C\tilde B_C^\top(\tilde B_T^{-1})^\top,
\]
故
\[
\det(I_{|V|-1}+XX^\top)
=\det(\tilde B_T\tilde B_T^\top)^{-1}\det(\tilde B_T\tilde B_T^\top+\tilde B_C\tilde B_C^\top).
\]
而 \(\det(\tilde B_T\tilde B_T^\top)=\det(\tilde B_T)^2=1\)，并且
\(
\tilde B_T\tilde B_T^\top+\tilde B_C\tilde B_C^\top=\tilde B\tilde B^\top
\)
为约化 Laplacian。
因此
\[
\operatorname{covol}(\Lambda(G))^2=\det(\tilde B\tilde B^\top).
\]
由 Kirchhoff 矩阵树定理，\(\det(\tilde B\tilde B^\top)=\tau(G)\)，从而
\(
\operatorname{covol}(\Lambda(G))=\sqrt{\tau(G)}
\)。

\textbf{（ii）Poisson 求和与极限）}
令 \(\Lambda\subset\RR^{r}\) 为任意满秩欧氏格，\(\Lambda^\ast\) 为其对偶格。
对高斯函数 \(\phi_t(x):=e^{-\pi t|x|^2}\) 有 Fourier 变换 \(\widehat\phi_t(\xi)=t^{-r/2}e^{-\pi|\xi|^2/t}\)。
应用 Poisson 求和公式得
\[
\Theta_\Lambda(t)=\sum_{z\in\Lambda}\phi_t(z)
=\frac{1}{\operatorname{covol}(\Lambda)}\sum_{y\in\Lambda^\ast}\widehat\phi_t(y)
=t^{-r/2}\operatorname{covol}(\Lambda)^{-1}\Theta_{\Lambda^\ast}(1/t).
\]
当 \(t\downarrow 0\) 时，\(1/t\uparrow\infty\)，并且
\[
\Theta_{\Lambda^\ast}(1/t)=1+O\bigl(e^{-\pi \lambda_1(\Lambda^\ast)^2/t}\bigr),
\]
其中 \(\lambda_1(\Lambda^\ast)\) 为对偶格最短非零向量长度。
取 \(c(G):=\pi\lambda_1(\Lambda(G)^\ast)^2\) 并代入 \(\operatorname{covol}(\Lambda(G))=\sqrt{\tau(G)}\) 即得所述两式。证毕。
\end{proof}

% Energy-shell lattice counting on affine cycle lattices; bit lower bound.

\begin{theorem}[回路圆维驱动的椭球格点计数：固定离散 Stokes 数据后的能量壳微观态熵]\label{thm:graph-energy-shell-lattice-counting}
令 \(G\) 为有限连通无自环多重图，并固定一组取向，使
\[
C_1(G;\ZZ)\cong \ZZ^{E},
\qquad
C_0(G;\ZZ)\cong \ZZ^{V},
\]
其中 \(E:=|E(G)|\)、\(V:=|V(G)|\)。
记边界算子（入射算子）为 \(\partial:C_1(G;\ZZ)\to C_0(G;\ZZ)\)，并记回路圆维
\[
r:=\operatorname{rank}_{\ZZ}H_1(G;\ZZ)=\operatorname{rank}_{\ZZ}\ker(\partial)=E-V+1.
\]
给定任意 \(\delta\in \mathrm{im}(\partial)\subset C_0(G;\ZZ)\)，任选一解 \(\varphi_0\in C_1(G;\ZZ)\) 使得 \(\partial\varphi_0=\delta\)。
取权重 \(w:E(G)\to(0,\infty)\)，并定义加权能量二次型
\[
\mathcal E_w(\varphi):=\sum_{e\in E(G)} w_e\,\varphi_e^2.
\]
考虑能量壳内的整数 \(1\)-形式计数
\[
N_{\delta,w}(E):=\card{\left\{\varphi\in C_1(G;\ZZ):\ \partial\varphi=\delta,\ \mathcal E_w(\varphi)\le E\right\}}.
\]
则存在正定矩阵 \(A\in\mathrm{Sym}_r^+(\RR)\)（仅依赖于 \(G\) 的回路格与权重 \(w\)，且与 \(\varphi_0\) 的选取无关）使得当 \(E\to\infty\) 时
\[
N_{\delta,w}(E)=\frac{\operatorname{Vol}(B_r)}{\sqrt{\det A}}\,E^{r/2}+O\!\left(E^{(r-1)/2}\right),
\]
其中 \(\operatorname{Vol}(B_r)=\pi^{r/2}/\Gamma(r/2+1)\) 为 \(\RR^r\) 中单位欧氏球体积。
并且若定义加权生成树多项式值
\[
\tau_{1/w}(G):=\sum_{T\ \mathrm{spanning\ tree}}\ \prod_{e\in T}w_e^{-1},
\]
则对任意由回路格 \(\ker(\partial)\subset \ZZ^{E}\) 的 \(\ZZ\)-基所诱导的 \(A\) 都有
\[
\boxed{\;\det A=\Bigl(\prod_{e\in E(G)}w_e\Bigr)\tau_{1/w}(G)\;}
\]
（定理 \ref{thm:graph-cycle-lattice-weighted-discriminant}，取 \(\rho=w\)）。
因此在固定 \(\delta\) 后，能量 \(E\) 内微观态熵满足
\[
\log N_{\delta,w}(E)=\frac r2\log E-\frac12\log\det A+O(1),
\]
其主阶系数仅由回路圆维 \(r\) 决定。
\end{theorem}
\begin{proof}
由于 \(\delta\in\mathrm{im}(\partial)\)，解集非空且为仿射格
\[
\{\varphi\in C_1(G;\ZZ):\ \partial\varphi=\delta\}=\varphi_0+\ker(\partial).
\]
取 \(\ker(\partial)\) 的一组 \(\ZZ\)-基并记相应基矩阵为 \(T\in\ZZ^{E\times r}\)，使得 \(\ker(\partial)=T\ZZ^r\)。
于是任意解可写为 \(\varphi=\varphi_0+Tk\)（\(k\in\ZZ^r\)）。
记 \(W:=\mathrm{diag}(w_e)\in\mathrm{Sym}_E^+(\RR)\)，则
\[
\mathcal E_w(\varphi_0+Tk)=(\varphi_0+Tk)^{\mathsf T}W(\varphi_0+Tk)
=k^{\mathsf T}(T^{\mathsf T}WT)k+2b^{\mathsf T}k+c,
\]
其中 \(A:=T^{\mathsf T}WT\in\mathrm{Sym}_r^+(\RR)\)、\(b:=T^{\mathsf T}W\varphi_0\)、\(c:=\varphi_0^{\mathsf T}W\varphi_0\)。
完成平方可将约束 \(\mathcal E_w(\varphi_0+Tk)\le E\) 改写为
\[
(k+\kappa)^{\mathsf T}A(k+\kappa)\le E',
\]
其中 \(\kappa\in\RR^r\) 与 \(E'=E+O(1)\) 由 \(A,b,c\) 决定。
因此 \(N_{\delta,w}(E)\) 等于一个以 \(A\) 定义的椭球与平移格 \(\ZZ^r+\kappa\) 的交点数。
对固定 \(A\) 的椭球族，Lipschitz 原理（体积近似）给出格点数等于椭球体积加上边界量级误差：
\[
N_{\delta,w}(E)=\operatorname{Vol}\{x\in\RR^r:\ x^{\mathsf T}Ax\le E'\}+O\!\left((E')^{(r-1)/2}\right).
\]
而体积可显式计算为
\[
\operatorname{Vol}\{x:\ x^{\mathsf T}Ax\le E'\}=\frac{\operatorname{Vol}(B_r)}{\sqrt{\det A}}\,(E')^{r/2}
=\frac{\operatorname{Vol}(B_r)}{\sqrt{\det A}}\,E^{r/2}+O\!\left(E^{(r-1)/2}\right).
\]
将误差并入即可得到主式。
由于 \(A=T^{\mathsf T}WT\) 仅依赖于回路格 \(\ker(\partial)\) 与权重 \(w\)，与特解 \(\varphi_0\) 无关。
对对数式，将主项取对数并吸收相对误差得到 \(\log N_{\delta,w}(E)=\frac r2\log E-\frac12\log\det A+O(1)\)。证毕。
\end{proof}
\leanverified{paper\_graph\_energy\_shell\_lattice\_counting}

\begin{corollary}[能量约束下的全息比特下界：圆维给出的有效自由度]\label{cor:graph-energy-shell-bit-lower-bound}
在定理 \ref{thm:graph-energy-shell-lattice-counting} 的设定下，任意能够区分集合
\[
\left\{\varphi\in C_1(G;\ZZ):\ \partial\varphi=\delta,\ \mathcal E_w(\varphi)\le E\right\}
\]
中全部元素的单射编码，其最小信息量（以比特计）满足
\[
\mathrm{bits}(E)\ \ge\ \log_2 N_{\delta,w}(E)=\frac r2\log_2 E+O(1).
\]
\end{corollary}
\begin{proof}
任意单射编码对 \(N_{\delta,w}(E)\) 个不同对象至少需要 \(\log_2 N_{\delta,w}(E)\) 比特。
再代入定理 \ref{thm:graph-energy-shell-lattice-counting} 的渐近式即得。证毕。
\end{proof}
\leanverified{paper\_graph\_energy\_shell\_bit\_lower\_bound}

\begin{proposition}[椭球约束下的 Gödel 码长极值：线性泛函的精确对偶]\label{prop:graph-energy-ellipsoid-godel-codelength-dual}
\leanverified{paper\_graph\_energy\_ellipsoid\_godel\_codelength\_dual}
在定理 \ref{thm:graph-energy-shell-lattice-counting} 的设定下，取回路格基坐标 \(k\in\ZZ^r\) 使得任意解可写为 \(\varphi=\varphi_0+Tk\)。
令 \(A=T^{\mathsf T}WT\in\mathrm{Sym}_r^+(\RR)\) 为能量二次型在回路坐标上的 Gram 矩阵。
取任意互异素数 \(p_1,\dots,p_r\)，并定义线性码长泛函
\[
\ell(k):=\sum_{j=1}^r k_j\log p_j=\langle k,\log p\rangle,
\qquad
\log p:=(\log p_1,\dots,\log p_r)\in\RR^r.
\]
则对任意 \(E>0\) 有精确极值公式
\[
\boxed{\;
\sup\{|\ell(x)|:\ x\in\RR^r,\ x^{\mathsf T}Ax\le E\}
=\sqrt E\,\sqrt{(\log p)^{\mathsf T}A^{-1}(\log p)}.
\;}
\]
从而对整数点集合
\(
\mathcal M_A(E):=\{k\in\ZZ^r:\ k^{\mathsf T}Ak\le E\}
\)
有支撑界
\[
|\ell(k)|\ \le\ \sqrt E\,\sqrt{(\log p)^{\mathsf T}A^{-1}(\log p)}
\qquad (k\in\mathcal M_A(E)).
\]
\end{proposition}
\begin{proof}
对任意满足 \(x^{\mathsf T}Ax\le E\) 的 \(x\in\RR^r\)，令 \(y:=A^{1/2}x\)。
则 \(|y|^2=x^{\mathsf T}Ax\le E\)，并且
\[
\ell(x)=\langle x,\log p\rangle=\langle y, A^{-1/2}\log p\rangle.
\]
由 Cauchy--Schwarz 不等式，
\[
|\ell(x)|
\le
|y|\,|A^{-1/2}\log p|
\le
\sqrt E\,\sqrt{(\log p)^{\mathsf T}A^{-1}(\log p)}.
\]
当取 \(y\) 与 \(A^{-1/2}\log p\) 共线且 \(|y|=\sqrt E\) 时达到上确界，故为精确极值。
对整数点支撑界由同一不等式直接推出。证毕。
\end{proof}

\begin{corollary}[行列式素支撑与核结构的不可约算术障碍]\label{cor:graph-cycle-lattice-determinant-prime-support}
在定理 \ref{thm:graph-energy-shell-lattice-counting} 的设定下，进一步假设 \(w_e\in\ZZ_{>0}\) 且所取回路格基使得 \(A=T^{\mathsf T}\mathrm{diag}(w_e)T\in M_r(\ZZ)\)。
考虑环面自同态
\[
F_A:\TT^r\to\TT^r,\qquad x\longmapsto Ax.
\]
则其核为有限阿贝尔群，且满足
\[
|\ker(F_A)|=|\det A|.
\]
并且对任意素数 \(\ell\)，\(\ker(F_A)\) 的 \(\ell\)-初等部分非平凡当且仅当 \(\ell\mid \det A\)。
因此 \(\det A\) 的素因子集合（等价地，由定理 \ref{thm:graph-cycle-lattice-weighted-discriminant}，
\(\det A=(\prod_e w_e)\tau_{1/w}(G)\) 的素因子集合）刻画了该核在素分解下的不可约算术支撑。
\end{corollary}
\begin{proof}
由 \(\TT^r=\RR^r/\ZZ^r\) 及 \(A\in M_r(\ZZ)\) 得 \(\ker(F_A)\cong \ZZ^r/A\ZZ^r\)。
对整数矩阵取 Smith 正规形可得 \(\ZZ^r/A\ZZ^r\) 为有限阿贝尔群，且其阶等于 \(|\det A|\)。
素分解刻画由有限阿贝尔群的标准分解定理得到：\(\ell\)-初等部分非平凡当且仅当该阶被 \(\ell\) 整除。
证毕。
\end{proof}

\leanverified{paper\_graph\_cycle\_lattice\_determinant\_prime\_support}

\begin{conjecture}[高圆维极限下的椭圆化观测：高斯--均匀解耦与“半经典”方差由有效阻抗决定]\label{conj:graph-energy-high-cdim-elliptic-observation}
考虑一族维数 \(r\to\infty\) 的正定矩阵 \(A_r\in\mathrm{Sym}_r^+(\RR)\) 与向量 \(a_r\in\RR^r\setminus\{0\}\)，并在能量预算 \(E_r\to\infty\) 下取均匀分布
\[
K_r\ \sim\ \mathrm{Unif}\bigl(\mathcal M_{A_r}(E_r)\bigr),
\qquad
\mathcal M_{A_r}(E_r):=\{k\in\ZZ^r:\ k^{\mathsf T}A_r k\le E_r\}.
\]
令线性观测（对数尺度）
\(
U_r:=\langle a_r,K_r\rangle
\)。
在“几何不退化”的尺度（例如 \(E_r/r\) 有界且远离 \(0\)，且 \((A_r)\) 的条件数在 \(r\) 上一致有界）下，预期成立如下极限律：
\begin{enumerate}
  \item \textbf{尺度边缘的中心极限定律.}
        归一化变量 \(\sqrt{r/E_r}\,U_r\) 收敛到中心高斯；其方差由对偶能量（有效阻抗）
        \(
        a_r^{\mathsf T}A_r^{-1}a_r
        \)
        决定。
        等价地，
        \(
        \mathrm{Var}(U_r/\sqrt{E_r})\asymp r^{-1}\,a_r^{\mathsf T}A_r^{-1}a_r
        \)。
  \item \textbf{相位的均匀化与渐近独立.}
        在将椭圆门输出以 Joukowsky 均匀化写作 \(z=\exp(U_r+\mathrm{i}\Theta_r)\)、\(w=z+z^{-1}\) 的口径下，
        相位变量 \(\Theta_r\) 的边缘分布趋于 \([0,2\pi)\) 上的均匀分布，且与 \(U_r\) 渐近独立。
  \item \textbf{椭圆侧极限分布.}
        因而 \(w\) 在椭圆族上的极限分布呈现“调和测度 \(\times\) 高斯径向涨落”的解耦结构：角向由调和测度给出，而尺度涨落由 \(a_r^{\mathsf T}A_r^{-1}a_r\) 控制。
\end{enumerate}
该极限律把回路圆维 \(r\) 解释为控制半经典噪声强度的主参量，并将“能量椭球--单标量读出--椭圆门外置”统一到可检验的概率型 bulk/boundary 制度之中。
\end{conjecture}

\endinput


\begin{remark}[离散面积律的圆维主项与生成树行列式校正]\label{rem:hypercube-boundary-cycle-entropy-area-law-tree-correction}
将定理 \ref{thm:graph-energy-shell-lattice-counting} 应用于边界邻接多图 \(G=\mathsf B_f\) 并固定散度数据（例如由体偏差或 Walsh 读出诱导的 \(\delta\)），可得能量阈值 \(E\) 内的可兼容边界整数 \(1\)-形式计数满足
\[
\log N_{\delta,w}(E)
=\frac{1}{2}\cdim\bigl(H_1(\mathsf B_f;\ZZ)\bigr)\log E
-\frac12\sum_{e\in E(\mathsf B_f)}\log w_e
-\frac12\log\tau_{1/w}(\mathsf B_f)
 +O(1),
\]
其中 \(\tau_{1/w}(\mathsf B_f)=\sum_{T}\prod_{e\in T}w_e^{-1}\) 为导通权 \(1/w\) 下的加权生成树多项式值（定理 \ref{thm:graph-cycle-lattice-weighted-discriminant}）。
该分解把“由边界回路圆维控制的幂律态数增长”与“由 Kirchhoff 复杂度给出的行列式修正”在同一可审计公式中对齐，并为正文关于面积律/纠缠语义接口的视界模型（例如 \S\ref{subsec:dephysicalized-horizon-timefiber-nohair} 的 AdS/CFT 解释性接口）提供一条纯离散的可计算原型。
\end{remark}

\begin{conjecture}[生成树行列式的连续谱极限接口]\label{conj:boundary-tree-determinant-spectral-limit}
设 \(\{\mathsf B^{(m)}_f\}_{m\ge 1}\) 为某一 dyadic 细化极限下的边界邻接多图序列，并取与边界归一化 Stokes 逆塔一致的权重归一化方案（参见附录 \ref{app:multiscale-normalized-stokes-inverse-tower}）。
记 \(L^{(m)}\) 为 \(\mathsf B^{(m)}_f\) 的加权 Laplacian（导通权为 \(1/w^{(m)}\)），并令 \(\tau_{1/w^{(m)}}(\mathsf B^{(m)}_f)=\det\bigl((L^{(m)})_{\mathrm{red}}\bigr)\) 为其加权生成树多项式值（Kirchhoff 定理）。
则在适当的重整化下，\(-\tfrac12\log\tau_{1/w^{(m)}}(\mathsf B^{(m)}_f)\) 收敛到某个极限量，该极限与连续边界算子（例如 Dirichlet-to-Neumann 映射或 Laplacian）的 \(\zeta\)-正则化行列式 \(\det_\zeta\) 相一致：
\[
-\frac12\log\tau_{1/w^{(m)}}(\mathsf B^{(m)}_f)\ \longrightarrow\ -\frac12\log\det_\zeta(\mathcal L_{\partial}).
\]
在该意义下，\(-\tfrac12\log\tau_{1/w}(\mathsf B_f)\) 可被视为对由圆维主项给出的“面积律”项的离散行列式修正。
\end{conjecture}

\begin{corollary}[回路格球截断的体积主项与 Gödel 外置对数规模]\label{cor:hypercube-cycle-lattice-counting-godel-log}
在定理 \ref{thm:hypercube-cycle-lattice-theta-asymptotic} 的设定下，对 \(R>0\) 定义球截断
\[
\Lambda(G;R):=\{z\in\Lambda(G):\ |z|\le R\}.
\]
则存在常数 \(C_G>0\) 使得当 \(R\to\infty\) 时有渐近计数
\[
\boxed{\;
|\Lambda(G;R)|
=\frac{\operatorname{vol}(B_r)}{\sqrt{\tau(G)}}\,R^r+O(R^{r-1}).
\;}
\]
其中 \(\operatorname{vol}(B_r)\) 为 \(\RR^{r}\) 中单位球体积。
进一步，取任意 Gödel 外置单射 \(\Gamma:\Lambda(G)\hookrightarrow\NN\)，则
\[
\boxed{\;
\log|\Gamma(\Lambda(G;R))|
=r\log R-\tfrac12\log\tau(G)+O(1).
\;}
\]
\leanverified{paper\_hypercube\_cycle\_lattice\_counting\_godel\_log}
\end{corollary}
\begin{proof}
记 \(\Lambda:=\Lambda(G)\subset \RR^{E}\)，其秩为 \(r\) 且 \(\operatorname{covol}(\Lambda)=\sqrt{\tau(G)}\)（定理 \ref{thm:hypercube-cycle-lattice-theta-asymptotic}）。
取 \(\Lambda\) 在其张成子空间上的一个基本平行体 \(P\)（可取由一组格基生成），并记其直径为 \(D\)。
令 \(B(R)\) 为该子空间中的半径 \(R\) 球。
\par
\textbf{（体积夹逼）}
由于 \(\{z+P:z\in\Lambda\}\) 构成无交叠的平移铺砌，有
\[
|\Lambda\cap B(R)|\cdot \operatorname{vol}(P)
\le \operatorname{vol}(B(R+D)).
\]
同理，若 \(z\in\Lambda\cap B(R-D)\)，则 \(z+P\subset B(R)\)，因此
\[
\operatorname{vol}(B(R-D))
\le |\Lambda\cap B(R)|\cdot \operatorname{vol}(P).
\]
合并并用 \(\operatorname{vol}(P)=\operatorname{covol}(\Lambda)\) 得
\[
\frac{\operatorname{vol}(B(R-D))}{\operatorname{covol}(\Lambda)}
\le |\Lambda(G;R)|
\le \frac{\operatorname{vol}(B(R+D))}{\operatorname{covol}(\Lambda)}.
\]
利用 \(\operatorname{vol}(B(R))=\operatorname{vol}(B_r)R^r\) 与二项展开，
\[
\operatorname{vol}(B(R\pm D))
=\operatorname{vol}(B_r)\,(R\pm D)^r
=\operatorname{vol}(B_r)R^r+O(R^{r-1}),
\]
从而得到首式，且隐含常数仅依赖于 \(G\)。
\par
\textbf{（Gödel 对数）}
由于 \(\Gamma\) 为单射，\(|\Gamma(\Lambda(G;R))|=|\Lambda(G;R)|\)。
对首式取对数即得
\[
\log|\Lambda(G;R)|
=\log\Bigl(\frac{\operatorname{vol}(B_r)}{\sqrt{\tau(G)}}R^r\bigl(1+O(R^{-1})\bigr)\Bigr)
=r\log R-\tfrac12\log\tau(G)+O(1).
\]
证毕。
\end{proof}

\begin{corollary}[Gödel 椭球截断的回路格点体积律]\label{cor:hypercube-cycle-lattice-godel-ellipsoid-counting}
\leanverified{paper\_hypercube\_cycle\_lattice\_godel\_ellipsoid\_counting}
令 \(G\) 为有限连通无自环多重图，并记
\[
r:=\operatorname{rank}_{\ZZ}H_1(G;\ZZ)=|E(G)|-|V(G)|+1.
\]
取互异素数 \(p_1,\dots,p_r\) 与指数向量 \(\mathbf r=(r_1,\dots,r_r)\in\NN^r\)，并定义 Gödel 整数
\[
N(\mathbf r):=\prod_{i=1}^{r}p_i^{r_i}.
\]
定义 Gödel 椭球
\[
E_r(\mathbf r;R):=\left\{x\in\RR^r:\ \sum_{i=1}^{r}\frac{x_i^2}{p_i^{r_i}}\le R^2\right\},
\]
并取任一欧氏等距同构 \(\iota:\RR^r\to \Span_{\RR}(\Lambda(G))\)（其中 \(\Lambda(G)\) 为定理 \ref{thm:hypercube-cycle-lattice-theta-asymptotic} 中的回路格），令
\(
E_G(\mathbf r;R):=\iota\bigl(E_r(\mathbf r;R)\bigr)
\)。
则存在常数 \(C_G>0\) 使得当 \(R\to\infty\) 时
\[
\boxed{\;
\card{\Lambda(G)\cap E_G(\mathbf r;R)}
=\frac{\operatorname{vol}(B_r)\sqrt{N(\mathbf r)}}{\sqrt{\tau(G)}}\,R^r
+O(R^{r-1}),\;}
\]
其中 \(\tau(G)\) 为生成树个数，\(\operatorname{vol}(B_r)\) 为 \(\RR^r\) 中单位球体积，且隐含常数仅依赖于 \(G\) 与所取欧氏结构。
从而
\[
\log\card{\Lambda(G)\cap E_G(\mathbf r;R)}
=r\log R+\tfrac12\log N(\mathbf r)-\tfrac12\log\tau(G)+\log\operatorname{vol}(B_r)+o(1).
\]
\par
特别地，取 \(G=\mathsf B_f\)（定义 \ref{def:hypercube-boundary-multigraph}）并记
\(
r(f):=\cdim(H_1(\mathsf B_f;\ZZ))
\)、\(
\tau_f:=\tau(\mathsf B_f)
\)，则上式在 \(r=r(f)\)、\(\tau(G)=\tau_f\) 下成立。
\end{corollary}

\begin{proof}
由定理 \ref{thm:hypercube-cycle-lattice-theta-asymptotic}，\(\Lambda(G)\) 的协体积为 \(\operatorname{covol}(\Lambda(G))=\sqrt{\tau(G)}\)。
取其一组格基生成的基本平行体 \(P\)，并记直径为 \(D\)。
对任意 \(R>D\)，体积夹逼（同推论 \ref{cor:hypercube-cycle-lattice-counting-godel-log} 的证明）给出
\[
\frac{\operatorname{vol}\bigl(E_r(\mathbf r;R-D)\bigr)}{\operatorname{covol}(\Lambda(G))}
\le
\card{\Lambda(G)\cap E_G(\mathbf r;R)}
\le
\frac{\operatorname{vol}\bigl(E_r(\mathbf r;R+D)\bigr)}{\operatorname{covol}(\Lambda(G))}.
\]
由于
\(
\operatorname{vol}(E_r(\mathbf r;R))=\operatorname{vol}(B_r)\sqrt{N(\mathbf r)}\,R^r
\)，
并且
\(
(R\pm D)^r=R^r+O(R^{r-1})
\)，
代入即得主式。
对主式取对数并吸收相对误差即得对数展开。证毕。
\end{proof}

\endinput


\begin{corollary}[Stokes torsion 的一回路计数律：直和回路格的球截断]\label{cor:hypercube-cycle-lattice-directsum-counting}
令 \(G\) 为有限连通多重图，并沿用定理 \ref{thm:hypercube-cycle-lattice-theta-asymptotic} 的回路格 \(\Lambda(G)\subset \RR^{E(G)}\)、秩 \(r\) 与生成树个数 \(\tau(G)\) 的记号。
对整数 \(m\ge 1\) 定义直和格
\[
\Lambda(G)^{\oplus m}:=\underbrace{\Lambda(G)\oplus\cdots\oplus\Lambda(G)}_{m\ \text{份}}
\subset \bigl(\RR^{E(G)}\bigr)^m,
\]
并以标准欧氏范数 \(|\cdot|_2\) 定义球截断
\[
\Lambda(G)^{\oplus m}(R):=\{h\in\Lambda(G)^{\oplus m}:\ |h|_2\le R\}.
\]
则当 \(R\to\infty\) 时有对数渐近式
\[
\boxed{\;
\log\bigl|\Lambda(G)^{\oplus m}(R)\bigr|
=mr\log R-\frac{m}{2}\log\tau(G)+O(1).\;}
\]
并且同一渐近式对任意仿射陪集 \(h_0+\Lambda(G)^{\oplus m}\) 的球截断计数同样成立。
\end{corollary}
\begin{proof}
由定理 \ref{thm:hypercube-cycle-lattice-theta-asymptotic}，\(\Lambda(G)\) 为秩 \(r\) 的欧氏格，且 \(\operatorname{covol}(\Lambda(G))=\sqrt{\tau(G)}\)。
直和格 \(\Lambda(G)^{\oplus m}\) 的秩为 \(mr\)，并且在正交直和的度量约定下满足协体积乘法性
\[
\operatorname{covol}\bigl(\Lambda(G)^{\oplus m}\bigr)=\operatorname{covol}(\Lambda(G))^{m}=\tau(G)^{m/2}.
\]
对格 \(\Lambda(G)^{\oplus m}\) 在其张成子空间中应用推论 \ref{cor:hypercube-cycle-lattice-counting-godel-log} 证明中的体积夹逼论证（将维数由 \(r\) 替换为 \(mr\)，协体积替换为 \(\tau(G)^{m/2}\)），得到
\[
\bigl|\Lambda(G)^{\oplus m}(R)\bigr|
=\frac{\operatorname{vol}(B_{mr})}{\tau(G)^{m/2}}\,R^{mr}+O(R^{mr-1}).
\]
取对数并吸收误差得所示对数渐近式。
对任意仿射陪集，平移不改变协体积与主导体积项，且截断边界层仅带来 \(O(R^{mr-1})\) 级别的修正，因此对数主项保持不变。证毕。
\end{proof}

\begin{corollary}[回路格截断的最小描述长度：树数对数扣除项]\label{cor:hypercube-cycle-lattice-directsum-mdl}
在推论 \ref{cor:hypercube-cycle-lattice-directsum-counting} 的设定下，取有限字母表 \(\mathcal A\)，\(\card{\mathcal A}=M\ge 2\)。
对任意 \(R\ge 2\)，若存在单射编码
\[
\mathrm{Enc}:\ \Lambda(G)^{\oplus m}(R)\hookrightarrow \mathcal A^{\le L},
\]
则必要长度满足
\[
L\ \ge\ \frac{\log\bigl|\Lambda(G)^{\oplus m}(R)\bigr|}{\log M}+O(1)
=\frac{mr\log R-\frac{m}{2}\log\tau(G)}{\log M}+O(1).
\]
\end{corollary}
\begin{proof}
单射性蕴含 \(\card{\mathcal A^{\le L}}\ge \bigl|\Lambda(G)^{\oplus m}(R)\bigr|\)。
另一方面，
\[
\card{\mathcal A^{\le L}}=\sum_{\ell=0}^{L}M^{\ell}=\frac{M^{L+1}-1}{M-1}\le \frac{M^{L+1}}{M-1},
\]
故 \(L\ge \frac{\log|\Lambda(G)^{\oplus m}(R)|}{\log M}+O(1)\)。
代入推论 \ref{cor:hypercube-cycle-lattice-directsum-counting} 的对数渐近式即得结论。证毕。
\end{proof}

\endinput



\begin{corollary}[回路圆维与规范体积的互补下界]\label{cor:hypercube-loop-cdim-gauge-volume-complement}
在命题 \ref{prop:hypercube-discrete-holographic-stokes-gauge-area} 与命题 \ref{prop:hypercube-boundary-multigraph-h1-cdim} 的设定下，有
\[
\boxed{\;
g(f)\ \ge\ (n\log 2-1)\ -\ 2\log 2\cdot\frac{\cdim(H_1(\mathsf B_f;\ZZ))+\card{X}-c(\mathsf B_f)}{2^n}\ +\ \frac{\card{X}}{2^n}.
\;}
\]
\end{corollary}
\leanverified{paper\_hypercube\_loop\_cdim\_gauge\_volume\_complement}

\begin{proof}
由命题 \ref{prop:hypercube-boundary-multigraph-h1-cdim} 得
\(
A(f)=\cdim(H_1(\mathsf B_f;\ZZ))+\card{X}-c(\mathsf B_f)
\)。
代入命题 \ref{prop:hypercube-discrete-holographic-stokes-gauge-area} 即得。证毕。
\end{proof}

\begin{corollary}[面积密度极限下的线性面积律]\label{cor:hypercube-linear-area-law-density}
取一列满射粗粒化 \(f_n:\Omega_n\twoheadrightarrow X_n\)。
记面积密度
\[
a_n:=\frac{A(f_n)}{n2^{n-1}}\in[0,1],
\]
以及每比特规范体积密度
\[
\bar g_n:=\frac{g(f_n)}{n}.
\]
若 \(\card{X_n}/2^n\to 0\)，则
\[
\boxed{\;
\liminf_{n\to\infty}\bar g_n\ \ge\ (1-\limsup_{n\to\infty}a_n)\log 2.
\;}
\]
特别地，若 \(a_n\to a\in[0,1]\)，则 \(\liminf_{n\to\infty}\bar g_n\ge (1-a)\log 2\)。
\end{corollary}
\leanverified{paper\_hypercube\_linear\_area\_law\_density}

\begin{proof}
由命题 \ref{prop:hypercube-discrete-holographic-stokes-gauge-area}，
\[
\bar g_n
\ge
\log 2-\frac{1}{n}-2\log 2\cdot\frac{A(f_n)}{n2^n}+\frac{\card{X_n}}{n2^n}.
\]
注意 \(\frac{A(f_n)}{2^n}=\frac{n}{2}a_n\)，故
\[
\bar g_n\ge (1-a_n)\log 2-\frac{1}{n}+\frac{\card{X_n}}{n2^n}.
\]
取 \(\liminf\) 并用 \(\card{X_n}/2^n\to 0\) 即得结论。证毕。
\end{proof}

\paragraph{Gödel--Stokes 体--界字典：通量偏差、Walsh 全息与回路扭量}
\label{par:fold-hypercube-godel-stokes-walsh-hodge}
本段在超立方体粗粒化的边界邻接多图 \(\mathsf B_f\) 上，给出一条完全离散的“边界 \(\to\) 体积”读出链：单坐标方向的有向通量严格等于体积偏差；Fourier--Walsh 系数可由边界有向 \(1\)-形式积分全息读出；固定体积偏差后边界层仍保留的自由度精确形成以 \(H_1(\mathsf B_f;\ZZ)\) 为平移群的扭量；并且上述数据可由素数寄存器作乘法外置，同时在加权能量下诱导显式椭球可行域。

\begin{proposition}[单坐标离散 Stokes 恒等式：有向通量等于体积偏差]\label{prop:fold-hypercube-godel-stokes-flux-bias}
\leanverified{paper\_fold\_hypercube\_godel\_stokes\_flux\_bias}
设 \(n\ge 1\) 且 \(\Omega_n:=\{0,1\}^n\)。
对 \(i\in\{1,\dots,n\}\)，记 \(\omega^{(i)}\) 为将 \(\omega\) 的第 \(i\) 位翻转后的顶点。
对任意 \(S\subseteq\Omega_n\)，定义第 \(i\) 坐标方向的有向边界计数
\[
N_i^+(S):=\#\{\omega\in S:\ \omega_i=0,\ \omega^{(i)}\notin S\},
\qquad
N_i^-(S):=\#\{\omega\in S:\ \omega_i=1,\ \omega^{(i)}\notin S\},
\]
并定义体积偏差
\[
b_i(S):=\#(S\cap\{\omega_i=0\})-\#(S\cap\{\omega_i=1\})
=\sum_{\omega\in S}(-1)^{\omega_i}\in\ZZ.
\]
则对一切 \(i\) 有严格恒等式
\[
\boxed{\;N_i^+(S)-N_i^-(S)=b_i(S)\;}
\]
从而无向边界数 \(c_i(S):=N_i^+(S)+N_i^-(S)\) 满足
\[
c_i(S)\ge |b_i(S)|.
\]
\end{proposition}
\begin{proof}
固定 \(i\)。映射 \(\{\omega:\omega_i=0\}\to\{\omega:\omega_i=1\}\)、\(\omega\mapsto\omega^{(i)}\) 为双射，故
\[
N_i^+(S)-N_i^-(S)
=\sum_{\omega:\omega_i=0}\bigl(\mathbf 1_S(\omega)-\mathbf 1_S(\omega^{(i)})\bigr)
=\#(S\cap\{\omega_i=0\})-\#(S\cap\{\omega_i=1\})
=b_i(S).
\]
再由 \(|a+b|\ge ||a|-|b||\)（取 \(a=N_i^+(S),b=N_i^-(S)\)）得 \(c_i(S)\ge |b_i(S)|\)。证毕。
\end{proof}

\leanverified{paper\_fold\_hypercube\_weighted\_discrete\_flux\_identity}
\begin{proposition}\label{prop:fold-hypercube-weighted-discrete-flux-identity}
设 \(n\ge 1\)，取参数 \(a=(a_1,\dots,a_n)\in(0,\infty)^n\)，并定义指数型权重
\[
w_a(\omega):=\prod_{j=1}^n a_j^{\omega_j},
\qquad
\omega\in\Omega_n.
\]
对任意 \(S\subseteq\Omega_n\) 与坐标 \(i\in\{1,\dots,n\}\)，定义加权定向边界通量
\[
N_i^+(S;a):=\sum_{\substack{\omega\in S\\ \omega_i=0,\ \omega^{(i)}\notin S}} w_a(\omega),
\qquad
N_i^-(S;a):=\sum_{\substack{\omega\in S\\ \omega_i=1,\ \omega^{(i)}\notin S}} w_a(\omega).
\]
则成立精确恒等式
\[
a_i\,N_i^+(S;a)-N_i^-(S;a)
=a_i\sum_{\substack{\omega\in S\\ \omega_i=0}}w_a(\omega)-\sum_{\substack{\omega\in S\\ \omega_i=1}}w_a(\omega).
\]
\end{proposition}
\begin{proof}
令 \(f=\mathbf 1_S\)，并考虑求和
\[
\Sigma:=\sum_{\omega:\ \omega_i=0} w_a(\omega)\bigl(f(\omega)-f(\omega^{(i)})\bigr).
\]
一方面按边界分解得
\[
\Sigma
=\sum_{\substack{\omega\in S\\ \omega_i=0,\ \omega^{(i)}\notin S}}w_a(\omega)
-\sum_{\substack{\omega\notin S\\ \omega_i=0,\ \omega^{(i)}\in S}}w_a(\omega)
=N_i^+(S;a)-\sum_{\substack{v\in S\\ v_i=1,\ v^{(i)}\notin S}}w_a(v^{(i)}).
\]
注意当 \(v_i=1\) 时 \(w_a(v^{(i)})=w_a(v)/a_i\)，故
\[
\Sigma=N_i^+(S;a)-\frac{1}{a_i}N_i^-(S;a).
\]
另一方面直接重排得
\[
\Sigma
=\sum_{\substack{\omega\in S\\ \omega_i=0}}w_a(\omega)
-\sum_{\substack{v\in S\\ v_i=1}}w_a(v^{(i)})
=\sum_{\substack{\omega\in S\\ \omega_i=0}}w_a(\omega)
-\frac{1}{a_i}\sum_{\substack{v\in S\\ v_i=1}}w_a(v).
\]
令两式相等并同乘 \(a_i\) 即得结论。证毕。
\end{proof}

\begin{corollary}[\texorpdfstring{$\ell^1$}{l1} 边界下界与坐标单调性的等号刻画]\label{cor:fold-hypercube-boundary-l1-bias-monotone}
在命题 \ref{prop:fold-hypercube-godel-stokes-flux-bias} 的记号下，记 \(\partial S\) 为超立方体图上 \(S\) 的无向边界边集（定义如 \eqref{eq:hypercube-gauge-volume-renyi-upperbound} 的证明中所用）。
则
\[
|\partial S|=\sum_{i=1}^n c_i(S)\ \ge\ \sum_{i=1}^n |b_i(S)|.
\]
并且对给定的 \(i\) 有
\[
c_i(S)=|b_i(S)|
\quad\Longleftrightarrow\quad
N_i^+(S)N_i^-(S)=0,
\]
亦即第 \(i\) 坐标方向的全部边界边在同一方向上外流。
若对所有 \(i\) 同时取等号，则存在坐标翻转自同构
\(
\omega_i\mapsto 1-\omega_i
\)
使得变换后的集合成为坐标偏序下的上闭集（逐坐标单调）。
\end{corollary}
\begin{proof}
第一条由逐坐标不等式 \(c_i(S)\ge |b_i(S)|\) 求和得到。
等号条件来自 \(|a+b|=||a|-|b||\iff ab=0\)（\(a,b\ge 0\)）。
当每个坐标方向的边界边均同向时，按坐标逐层检验可得：在适当坐标翻转后，沿任一坐标由 \(0\) 到 \(1\) 的邻接翻转不会离开集合，从而集合为上闭。证毕。
\end{proof}
\leanverified{paper\_fold\_hypercube\_boundary\_l1\_bias\_monotone}

\begin{proposition}[Gödel 边界通量码与乘法体积读出]\label{prop:fold-hypercube-godel-flux-multiplicative-volume}
\leanverified{paper\_fold\_hypercube\_godel\_flux\_multiplicative\_volume}
取一组互异素数 \((p_i)_{i=1}^n\)，并定义 Gödel 边界通量码
\[
\Phi_p(S):=\prod_{i=1}^n p_i^{\,N_i^+(S)-N_i^-(S)}\in\QQ_{>0}.
\]
则 \(\Phi_p(S)=\prod_{i=1}^n p_i^{\,b_i(S)}\)，因而仅由体积偏差向量 \((b_1(S),\dots,b_n(S))\) 决定。
进一步定义顶点 Gödel 编码
\[
G(\omega):=\prod_{i=1}^n p_i^{\omega_i}\in\NN_{>0},
\]
则有严格恒等式
\[
\boxed{\;
\prod_{\omega\in S}G(\omega)
=\prod_{i=1}^n p_i^{\,\#(S\cap\{\omega_i=1\})}
=\prod_{i=1}^n p_i^{\frac{|S|-b_i(S)}{2}}\in\NN_{>0}.
\;}
\]
因此 \(\prod_{\omega\in S}G(\omega)\) 被 \(|S|\) 与 \(\Phi_p(S)\) 等价确定（在对数坐标上为一条线性关系）。
\end{proposition}
\begin{proof}
由命题 \ref{prop:fold-hypercube-godel-stokes-flux-bias} 得 \(N_i^+-N_i^-=b_i\)，代回即得 \(\Phi_p\) 的表达式。
另一方面，
\[
\prod_{\omega\in S}G(\omega)
=\prod_{\omega\in S}\prod_{i=1}^n p_i^{\omega_i}
=\prod_{i=1}^n p_i^{\sum_{\omega\in S}\omega_i}
=\prod_{i=1}^n p_i^{\#(S\cap\{\omega_i=1\})}.
\]
又由
\(
b_i(S)=\#(S\cap\{\omega_i=0\})-\#(S\cap\{\omega_i=1\})=|S|-2\#(S\cap\{\omega_i=1\})
\)
得 \(\#(S\cap\{\omega_i=1\})=(|S|-b_i(S))/2\)，代回即得结论。证毕。
\end{proof}

\begin{proposition}[体积约束诱导的偏差二次约束与 Gödel 码长椭球上界]\label{prop:fold-hypercube-bias-ellipsoid-godel-length}
设 \(S\subseteq\Omega_n\) 且 \(|S|=s\)，并仍记体积偏差 \(b_i(S)=\sum_{\omega\in S}(-1)^{\omega_i}\)（命题 \ref{prop:fold-hypercube-godel-stokes-flux-bias}）。
则成立严格二次型约束
\[
\boxed{\;\sum_{i=1}^n b_i(S)^2\ \le\ s(2^n-s).\;}
\]
进一步，对任意互异素数族 \((p_i)_{i=1}^n\) 定义 Gödel 外置码长
\[
\mathcal{L}_{\mathrm G}(S)
:=\log\Bigl(\prod_{i=1}^n p_i^{|b_i(S)|}\Bigr)
=\sum_{i=1}^n |b_i(S)|\log p_i,
\]
则有椭球型上界
\[
\boxed{\;\mathcal{L}_{\mathrm G}(S)\ \le\
\Bigl(\sum_{i=1}^n(\log p_i)^2\Bigr)^{1/2}\cdot \sqrt{s(2^n-s)}.\;}
\]
\end{proposition}
\begin{proof}
令 \(f:=\mathbf 1_S\)，并在 \(\Omega_n\) 上取均匀期望 \(\EE\)。
则对单点集合 \(\{i\}\) 的 Fourier--Walsh 系数有
\[
\widehat f(\{i\})
=\EE\bigl[f(\omega)(-1)^{\omega_i}\bigr]
=2^{-n}\sum_{\omega\in S}(-1)^{\omega_i}
=2^{-n}b_i(S).
\]
Parseval 恒等式给出
\[
\sum_{I\subseteq[n]}\widehat f(I)^2=\EE[f^2]=\EE[f]=\frac{s}{2^n}=:p,
\qquad
\widehat f(\varnothing)=p,
\]
从而
\[
\sum_{I\neq\varnothing}\widehat f(I)^2=p-p^2=p(1-p).
\]
特别地，
\[
\sum_{i=1}^n \widehat f(\{i\})^2\le p(1-p).
\]
两边乘以 \(2^{2n}\) 得
\[
\sum_{i=1}^n b_i(S)^2\le 2^{2n}p(1-p)=s(2^n-s).
\]
对码长，用 Cauchy--Schwarz 不等式
\[
\mathcal{L}_{\mathrm G}(S)
=\sum_i |b_i(S)|\log p_i
\le
\Bigl(\sum_i b_i(S)^2\Bigr)^{1/2}\Bigl(\sum_i(\log p_i)^2\Bigr)^{1/2},
\]
并代入前式即得。证毕。
\end{proof}

\leanverified{paper\_fold\_hypercube\_bias\_ellipsoid\_godel\_length}

\begin{corollary}[粗粒化切割面积的偏差下界]\label{cor:fold-hypercube-cutedges-bias-l1}
\leanverified{paper\_fold\_hypercube\_cutedges\_bias\_l1}
沿用定义 \ref{def:hypercube-coarsegraining-area-gauge-volume} 的满射粗粒化 \(f:\Omega_n\twoheadrightarrow X\)，并记纤维 \(S_x=f^{-1}(x)\)。
则
\[
2A(f)=\sum_{x\in X}|\partial S_x|
\ \ge\
\sum_{x\in X}\sum_{i=1}^n |b_i(S_x)|.
\]
\end{corollary}
\begin{proof}
对每个纤维 \(S_x\) 用推论 \ref{cor:fold-hypercube-boundary-l1-bias-monotone} 的下界 \(|\partial S_x|\ge \sum_i |b_i(S_x)|\)。
再注意每条切割边恰在两条纤维边界中各计一次，故 \(\sum_x|\partial S_x|=2A(f)\)。证毕。
\end{proof}

\begin{theorem}[加权 Walsh 模式的广义 Stokes 公式：Fourier 系数的边界全息读出]\label{thm:fold-hypercube-weighted-walsh-stokes}
\leanverified{paper\_fold\_hypercube\_weighted\_walsh\_stokes}
在超立方体上赋予坐标边权 \(w_i>0\)，并定义加权拉普拉斯算子
\[
(\Delta_w\psi)(\omega):=\sum_{i=1}^n w_i\bigl(\psi(\omega)-\psi(\omega^{(i)})\bigr).
\]
对非空指标集 \(I\subseteq\{1,\dots,n\}\)，令 Walsh 字符
\[
\chi_I(\omega):=(-1)^{\sum_{i\in I}\omega_i}.
\]
则 \(\chi_I\) 为 \(\Delta_w\) 的本征函数，且
\[
\Delta_w\chi_I=2\Bigl(\sum_{i\in I}w_i\Bigr)\chi_I.
\]
对任意 \(S\subseteq\Omega_n\)，令 \(\partial^{\to}S\) 表示从 \(S\) 指向补集的外法向有向边集合。
定义边界上的有向 \(1\)-形式
\[
\alpha_{I,w}(\omega\to\omega^{(i)})
:=\frac{w_i}{2\sum_{j\in I}w_j}\bigl(\chi_I(\omega^{(i)})-\chi_I(\omega)\bigr),
\qquad (\omega\to\omega^{(i)})\in\partial^{\to}S.
\]
则成立广义离散 Stokes 公式
\[
\boxed{\;\sum_{\omega\in S}\chi_I(\omega)=\sum_{e\in \partial^{\to}S}\alpha_{I,w}(e).\;}
\]
特别地，对 \(f=\mathbf 1_S\) 的 Fourier--Walsh 系数
\(
\widehat f(I)=2^{-n}\sum_{\omega\in\Omega_n}f(\omega)\chi_I(\omega)
\)
有边界积分表达
\[
2^n\widehat{\mathbf 1_S}(I)=\sum_{e\in \partial^{\to}S}\alpha_{I,w}(e).
\]
\end{theorem}
\begin{proof}
由本征关系，
\[
2\Bigl(\sum_{i\in I}w_i\Bigr)\sum_{\omega\in S}\chi_I(\omega)
=\sum_{\omega\in S}(\Delta_w\chi_I)(\omega)
=\sum_{\omega\in S}\sum_{i=1}^n w_i\bigl(\chi_I(\omega)-\chi_I(\omega^{(i)})\bigr).
\]
将内层求和按无向边配对，内部边贡献相消，仅剩跨越 \(S\) 与补集的边；并按外向取向把每条割边写成 \((\omega\to\omega^{(i)})\in\partial^{\to}S\)，得到
\[
\sum_{\omega\in S}(\Delta_w\chi_I)(\omega)
=\sum_{(\omega\to\omega^{(i)})\in\partial^{\to}S}w_i\bigl(\chi_I(\omega)-\chi_I(\omega^{(i)})\bigr).
\]
两边同除 \(2\sum_{i\in I}w_i\) 即得所述公式（与 \(\alpha_{I,w}\) 仅差号，已由定义吸收）。证毕。
\end{proof}

% Green inversion on the weighted hypercube from boundary currents.

\begin{theorem}[加权超立方体的 Green 反演：由边界电流到点值的唯一重构]\label{thm:fold-hypercube-green-inversion-unique-reconstruction}
\leanverified{paper\_fold\_hypercube\_green\_inversion\_unique\_reconstruction}
设 \(n\ge 1\) 且 \(\Omega_n:=\{0,1\}^n\)。
取权重 \(w_1,\dots,w_n>0\)，并定义加权拉普拉斯算子
\[
(\Delta_w u)(\omega):=\sum_{i=1}^n w_i\bigl(u(\omega)-u(\omega^{(i)})\bigr),
\qquad \omega\in\Omega_n,
\]
其中 \(\omega^{(i)}\) 表示翻转第 \(i\) 位后的顶点。
对非空指标集 \(I\subseteq[n]\) 定义 Walsh 字符
\[
\chi_I(\omega):=(-1)^{\sum_{i\in I}\omega_i},
\]
则 \(\chi_I\) 为 \(\Delta_w\) 的本征函数且满足
\[
\Delta_w\chi_I=2\Bigl(\sum_{i\in I}w_i\Bigr)\chi_I.
\]
令 Green 核（在常数模态上取零）为
\[
G_w(\omega,\eta):=2^{-n}\sum_{\varnothing\ne I\subseteq[n]}
\frac{\chi_I(\omega)\chi_I(\eta)}{2\sum_{j\in I}w_j},
\qquad \omega,\eta\in\Omega_n.
\]
则对任意 \(f:\Omega_n\to\RR\) 与任意锚点 \(\omega_\ast\in\Omega_n\)，若定义有向边电流
\[
J_f(\omega\to\omega^{(i)}):=f(\omega)-f(\omega^{(i)}),
\qquad (\omega\to\omega^{(i)})\ \text{为超立方体有向边},
\]
则成立精确反演公式
\[
f(\omega)
=f(\omega_\ast)
+\sum_{i=1}^n \sum_{\substack{u\in\Omega_n\\ u_i=0}}
w_i\,J_f(u\to u^{(i)})
\Bigl(
G_w(\omega,u)-G_w(\omega,u^{(i)})-G_w(\omega_\ast,u)+G_w(\omega_\ast,u^{(i)})
\Bigr).
\]
从而，给定全边电流 \(J_f\) 与单点值 \(f(\omega_\ast)\)，函数 \(f\) 被唯一确定。
\end{theorem}
\begin{proof}
对任意 \(u:\Omega_n\to\RR\) 作 Walsh 展开 \(u=\widehat u(\varnothing)+\sum_{I\ne\varnothing}\widehat u(I)\chi_I\)。
由本征关系可得对每个 \(I\ne\varnothing\),
\[
\widehat{\Delta_w u}(I)
=2\Bigl(\sum_{i\in I}w_i\Bigr)\widehat u(I),
\]
并且 \(\widehat{\Delta_w u}(\varnothing)=0\)。
由 Green 核的定义，直接计算得到对每个 \(I\ne\varnothing\),
\[
\widehat{G_w(\omega,\cdot)}(I)=\frac{2^{-n}\chi_I(\omega)}{2\sum_{i\in I}w_i},
\]
从而在 \(\eta\) 变量上有恒等式
\[
u(\omega)-u(\omega_\ast)
=\sum_{\eta\in\Omega_n}\bigl(G_w(\omega,\eta)-G_w(\omega_\ast,\eta)\bigr)\,(\Delta_w u)(\eta).
\]
将 \((\Delta_w u)(\eta)=\sum_i w_i(u(\eta)-u(\eta^{(i)}))\) 代入，并把求和按无向边 \(\{u,u^{(i)}\}\) 配对整理，可得
\[
u(\omega)-u(\omega_\ast)
=\sum_{i=1}^n \sum_{\substack{u\in\Omega_n\\ u_i=0}}
w_i\,(u(u)-u(u^{(i)}))
\Bigl(
G_w(\omega,u)-G_w(\omega,u^{(i)})-G_w(\omega_\ast,u)+G_w(\omega_\ast,u^{(i)})
\Bigr).
\]
取 \(u=f\) 即得所示公式；唯一性由右端仅依赖于 \(J_f\) 与锚点值立得。证毕。
\end{proof}

\endinput


\begin{theorem}[全 Walsh 纤维谱的边界可微性：谱向量落入切空间并诱导回路扭量]\label{thm:fold-hypercube-fiber-walsh-divergence-torsor}
\leanverified{paper\_fold\_hypercube\_fiber\_walsh\_divergence\_torsor}
沿用定义 \ref{def:hypercube-coarsegraining-area-gauge-volume} 的满射粗粒化 \(f:\Omega_n\twoheadrightarrow X\) 及其边界邻接多图 \(\mathsf B_f\)（定义 \ref{def:hypercube-boundary-multigraph}）。
对任意非空 \(I\subseteq[n]\)，定义纤维 Walsh 谱向量
\[
C_I\in \RR^{X},
\qquad
C_I(x):=\sum_{\omega\in S_x}\chi_I(\omega),
\]
其中 \(S_x=f^{-1}(x)\)。
将 \(C_0(\mathsf B_f;\RR)\) 与 \(\RR^{X}\) 视为同一对象，并取边界算子 \(\partial:C_1(\mathsf B_f;\RR)\to C_0(\mathsf B_f;\RR)\)。
则：
\begin{enumerate}
  \item \label{it:fiber-walsh-div-existence} 存在 \(F_I\in C_1(\mathsf B_f;\RR)\) 使得
  \[
  \boxed{\;\partial F_I=C_I.\;}
  \]
  并且 \(F_I\) 可由定理 \ref{thm:fold-hypercube-weighted-walsh-stokes} 的边界 \(1\)-形式 \(\alpha_{I,w}\) 通过将超立方体切割边逐条推前到 \(\mathsf B_f\) 的有向多重边并加总得到。
  \item \label{it:fiber-walsh-torsor} 解集为仿射空间
  \[
  \{F\in C_1(\mathsf B_f;\RR):\ \partial F=C_I\}
  =F_I+\ker(\partial),
  \qquad
  \ker(\partial)\cong H_1(\mathsf B_f;\RR).
  \]
  \item \label{it:fiber-walsh-dim-collapse} 维数塌缩：记 \(c(\mathsf B_f)\) 为连通分量数，则
  \[
  \boxed{\;\dim_{\RR}\operatorname{span}\{C_I:\ I\neq\varnothing\}\le |X|-c(\mathsf B_f).\;}
  \]
\end{enumerate}
\end{theorem}
\begin{proof}
\textbf{（\ref{it:fiber-walsh-div-existence}）}
固定 \(I\neq\varnothing\) 并任取 \(\mathsf B_f\) 的一组取向，使每条无向多重边被定向为某个有向边 \(e=(x\to y)\)。
对该有向边，令 \(E_{x\to y}\) 为所有从纤维 \(S_x\) 指向纤维 \(S_y\) 的超立方体有向割边集合：
\[
E_{x\to y}:=\{(\omega\to\omega^{(i)}):\ f(\omega)=x,\ f(\omega^{(i)})=y\}.
\]
在 \(C_1(\mathsf B_f;\RR)\) 上定义 \(F_I\) 的边系数为
\[
F_I(x\to y):=\sum_{e\in E_{x\to y}}\alpha_{I,w}(e),
\]
其中 \(\alpha_{I,w}\) 取自定理 \ref{thm:fold-hypercube-weighted-walsh-stokes}，并按多重边逐条加总。
对任意 \(x\in X\)，由边界算子定义，
\[
(\partial F_I)(x)
=\sum_{y}F_I(x\to y)-\sum_{y}F_I(y\to x)
=\sum_{e\in\partial^{\to}S_x}\alpha_{I,w}(e),
\]
其中最后一步利用：对跨越 \(S_x\) 与补集的每条无向割边，其在第二项中以相反取向出现；而 \(\alpha_{I,w}\) 对反向边取相反值，故两项合并恰为 \(S_x\) 的外法向边界求和。
由定理 \ref{thm:fold-hypercube-weighted-walsh-stokes} 取 \(S=S_x\) 即得
\(
(\partial F_I)(x)=\sum_{\omega\in S_x}\chi_I(\omega)=C_I(x)
\)。
由于 \(x\) 任意，得 \(\partial F_I=C_I\)。

\textbf{（\ref{it:fiber-walsh-torsor}）}
若 \(\partial F=\partial F'\)，则 \(\partial(F-F')=0\)，故差落入 \(\ker(\partial)\)；
反之对任意 \(Z\in\ker(\partial)\) 有 \(\partial(F_I+Z)=C_I\)，从而解集为 \(F_I+\ker(\partial)\)。
图作为一维 CW 复形满足 \(\mathrm{im}(\partial_2)=0\)，故 \(\ker(\partial)\cong H_1(\mathsf B_f;\RR)\)。

\textbf{（\ref{it:fiber-walsh-dim-collapse}）}
由 \(\partial F_I=C_I\) 得 \(C_I\in\mathrm{im}(\partial)\)。
而对任意一维 CW 复形，有
\(
\dim_{\RR}\mathrm{im}(\partial)=|X|-c(\mathsf B_f)
\)；
因此 \(\operatorname{span}\{C_I\}\subseteq \mathrm{im}(\partial)\) 给出所述维数上界。证毕。
\end{proof}

\begin{proposition}[Gödel--Hodge 扭量：固定偏差后的边界自由度由回路群承载]\label{prop:fold-hypercube-godel-hodge-torsor}
\leanverified{paper\_fold\_hypercube\_godel\_hodge\_torsor}
沿用定义 \ref{def:hypercube-boundary-multigraph} 的边界邻接多图 \(\mathsf B_f\)，并取其一维链群 \(C_1(\mathsf B_f;\ZZ)\) 与边界算子 \(\partial:C_1\to C_0\)。
对每个坐标 \(i\)，定义纤维偏差向量
\[
b^{(i)}\in C_0(\mathsf B_f;\ZZ),
\qquad
b^{(i)}_x:=b_i(S_x).
\]
则存在一个自然整数流 \(F^{(i)}\in C_1(\mathsf B_f;\ZZ)\)（由超立方体的第 \(i\) 坐标割边按 \(0\to 1\) 取向并推前到 \(\mathsf B_f\) 得到），满足
\[
\partial F^{(i)}=b^{(i)}.
\]
并且若解集
\[
\mathcal S_i:=\{F\in C_1(\mathsf B_f;\ZZ):\ \partial F=b^{(i)}\}
\]
非空，则其为仿射格（torsor）
\[
\mathcal S_i=F^{(i)}+\ker\partial,
\qquad
\ker\partial\cong H_1(\mathsf B_f;\ZZ).
\]
从而对全部坐标同时固定 \((b^{(1)},\dots,b^{(n)})\) 时，残留的全局边界自由度构成以 \(H_1(\mathsf B_f;\ZZ)^n\) 为平移群的扭量；其自由秩为
\[
\operatorname{rank}_{\ZZ}H_1(\mathsf B_f;\ZZ)^n
=n\bigl(A(f)-|X|+c(\mathsf B_f)\bigr)
\]
（命题 \ref{prop:hypercube-boundary-multigraph-h1-cdim}）。
\end{proposition}
\begin{proof}
把超立方体的第 \(i\) 坐标割边按 \(0\to 1\) 方向计为有向 \(1\)-链，并通过端点纤维标记 \(f:\Omega_n\to X\) 推前到 \(\mathsf B_f\) 的边链，得 \(F^{(i)}\)。
对任一顶点 \(x\in X\)，\((\partial F^{(i)})_x\) 等于从纤维 \(S_x\) 发出的 \(0\to 1\) 边数减去 \(1\to 0\) 边数，按命题 \ref{prop:fold-hypercube-godel-stokes-flux-bias} 恰为 \(b_i(S_x)=b_x^{(i)}\)。
若 \(F,F'\) 同为解，则 \(\partial(F-F')=0\)，故差落入 \(\ker\partial\)；反之加上任意 \(Z\in\ker\partial\) 仍为解，得到扭量结构。
最后 \(\ker\partial\cong H_1(\mathsf B_f;\ZZ)\) 与秩公式由命题 \ref{prop:hypercube-boundary-multigraph-h1-cdim} 给出。证毕。
\end{proof}

\begin{theorem}[边界邻接多图的有效电阻范数与最小能量原理]\label{thm:fold-hypercube-boundary-multigraph-effective-resistance-min-energy}
\leanverified{paper\_fold\_hypercube\_boundary\_multigraph\_effective\_resistance\_min\_energy}
设 \(G=(V,E)\) 为连通有限无向多图，取一组取向将其视作有向边集 \(E\)。
记边界算子 \(\partial:\RR^{E}\to \RR^{V}\) 为对应的入射（incidence）线性算子。
给定导纳（conductance）\(c_e>0\)（\(e\in E\)），令 \(C:=\mathrm{diag}(c_e)_{e\in E}\)，并定义加权 Laplace 算子
\[
L_c:=\partial\,C\,\partial^{\mathsf T}:\RR^{V}\to \RR^{V}.
\]
对满足 \(\sum_{v\in V}\delta(v)=0\) 的散度数据 \(\delta\in\RR^{V}\)，定义有效电阻范数
\[
\|\delta\|_{\mathrm{eff}}^{2}:=\delta^{\mathsf T}L_c^{+}\delta,
\]
其中 \(L_c^{+}\) 为 \(L_c\) 的 Moore--Penrose 伪逆。
则在所有满足 \(\partial\phi=\delta\) 的实值流 \(\phi\in\RR^{E}\) 中，二次能量
\[
\mathcal E_c(\phi):=\sum_{e\in E}\frac{\phi(e)^2}{c_e}
\]
存在唯一极小值，并满足
\[
\boxed{\;\min_{\partial\phi=\delta}\mathcal E_c(\phi)=\|\delta\|_{\mathrm{eff}}^{2}.\;}
\]
极小流必为梯度流：存在势函数 \(u\in\RR^{V}\) 使得
\[
\phi^\star=C\,\partial^{\mathsf T}u,
\qquad
L_cu=\delta,
\]
其中 \(u\) 唯一到加常数。
特别地，对边界邻接多图 \(G=\mathsf B_f\) 取 \(c_e\) 为对应无向多重边的重数，并对每个坐标 \(i\) 取命题 \ref{prop:fold-hypercube-godel-hodge-torsor} 的偏差向量 \(b^{(i)}\) 与整数流 \(F^{(i)}\)，则
\[
\|b^{(i)}\|_{\mathrm{eff}}^{2}\ \le\ \mathcal E_c(F^{(i)}).
\]
若进一步对每条边 \(e\) 有 \(0\le F^{(i)}(e)\le c_e\)（例如 \(F^{(i)}\) 由超立方体第 \(i\) 坐标割边推前得到时成立），则还满足
\[
\mathcal E_c(F^{(i)})\le \sum_{e\in E}F^{(i)}(e).
\]
\end{theorem}
\begin{proof}
约束集 \(\{\phi\in\RR^{E}:\ \partial\phi=\delta\}\) 为仿射子空间，\(\mathcal E_c\) 为严格凸二次型，故极小解存在且唯一。
对拉格朗日函数
\(
\mathcal L(\phi,u)=\sum_e \phi(e)^2/c_e+2\langle u,\delta-\partial\phi\rangle
\)
取一阶条件得
\(
2C^{-1}\phi-2\partial^{\mathsf T}u=0
\)，即 \(\phi=C\partial^{\mathsf T}u\)。
代入约束得 \(L_cu=\partial C\partial^{\mathsf T}u=\delta\)，解唯一到常数（因 \(G\) 连通且 \(\ker L_c=\langle \mathbf 1\rangle\)）。
将 \(\phi^\star=C\partial^{\mathsf T}u\) 代回能量并用伪逆表述可得
\[
\mathcal E_c(\phi^\star)=u^{\mathsf T}\delta=\delta^{\mathsf T}L_c^{+}\delta=\|\delta\|_{\mathrm{eff}}^{2}.
\]
对特殊化 \((\delta,\phi)=(b^{(i)},F^{(i)})\)，由极小性立即得 \(\|b^{(i)}\|_{\mathrm{eff}}^{2}\le \mathcal E_c(F^{(i)})\)。
若对每条边有 \(0\le F^{(i)}(e)\le c_e\)，则 \(F^{(i)}(e)^2/c_e\le F^{(i)}(e)\)，逐边求和得到 \(\mathcal E_c(F^{(i)})\le \sum_e F^{(i)}(e)\)。证毕。
\end{proof}

\begin{proposition}[回路自由度的素数预算下界：自由秩不可压缩]\label{prop:fold-hypercube-h1-prime-budget-minimal}
\leanverified{paper\_fold\_hypercube\_h1\_prime\_budget\_minimal}
沿用命题 \ref{prop:fold-hypercube-godel-hodge-torsor} 记号，并令
\[
r:=\operatorname{rank}_{\ZZ}H_1(\mathsf B_f;\ZZ)=A(f)-|X|+c(\mathsf B_f).
\]
则 \(H_1(\mathsf B_f;\ZZ)^n\cong \ZZ^{nr}\)。
取互异素数族 \(\{q_{i,j}\}_{1\le i\le n,\,1\le j\le r}\)，并定义映射
\[
\Gamma:\ \ZZ^{nr}\longrightarrow \QQ_{>0},
\qquad
(a_{i,j})\longmapsto \prod_{i=1}^n\prod_{j=1}^r q_{i,j}^{a_{i,j}}.
\]
则 \(\Gamma\) 为单射群同态。
并且若 \(A\) 为任意自由阿贝尔群且存在单射 \(\ZZ^{nr}\hookrightarrow A\)，则必有 \(\operatorname{rank}_{\ZZ}A\ge nr\)；
等价地，以互异素数作为独立乘法通道对 \(H_1(\mathsf B_f;\ZZ)^n\) 作 Gödel 外置时，所需通道数下界恰为 \(nr\)。
\end{proposition}
\begin{proof}
由唯一分解定理，\(\langle q_{i,j}\rangle\subset\QQ_{>0}\) 的自由部分同构于 \(\ZZ^{nr}\)，故 \(\Gamma\) 单射。
对最小性，任意单射 \(B\hookrightarrow A\)（\(A,B\) 自由阿贝尔）在 \(\QQ\)-张量下诱导线性单射 \(B\otimes\QQ\hookrightarrow A\otimes\QQ\)，故 \(\dim_{\QQ}(B\otimes\QQ)\le \dim_{\QQ}(A\otimes\QQ)\)，即 \(\operatorname{rank}_{\ZZ}A\ge \operatorname{rank}_{\ZZ}B=nr\)。证毕。
\end{proof}

\begin{proposition}[加权边界能量诱导的椭球约束]\label{prop:fold-hypercube-weighted-energy-ellipsoid-bias}
对权重 \(w_i>0\) 定义加权边界能量
\[
E_w(S):=\sum_{i=1}^n w_i\,c_i(S),
\]
其中 \(c_i(S)\) 为第 \(i\) 坐标方向无向边界数（命题 \ref{prop:fold-hypercube-godel-stokes-flux-bias}）。
则成立二次型下界
\[
\boxed{\;
E_w(S)\ \ge\ 2^{1-n}\sum_{i=1}^n w_i\,b_i(S)^2.
\;}
\]
从而在给定 \(E_w(S)\) 时，偏差向量 \(b(S)\) 必落入轴对齐椭球
\[
\sum_{i=1}^n w_i\,b_i^2\ \le\ 2^{n-1}E_w(S).
\]
\end{proposition}
\begin{proof}
令 \(f=\mathbf 1_S\)。由于 \(f(\omega)-f(\omega^{(i)})\in\{0,\pm 1\}\)，有
\[
2E_w(S)=\sum_{i=1}^n w_i\sum_{\omega\in\Omega_n}\bigl(f(\omega)-f(\omega^{(i)})\bigr)^2.
\]
将其按 \(2^{-n}\) 归一化得到加权 Dirichlet 形式
\(
\mathcal E_w(f)=2^{-n}\cdot 2E_w(S)=2^{1-n}E_w(S)
\)。
在 Walsh 基底上（参见定理 \ref{thm:discussion-hypercube-stokes-fourier-binomial} 的谱层展开口径），有
\[
\mathcal E_w(f)=\sum_{\varnothing\ne I\subseteq[n]}4\Bigl(\sum_{i\in I}w_i\Bigr)\,|\widehat f(I)|^2
\ \ge\
\sum_{i=1}^n 4w_i\,|\widehat f(\{i\})|^2.
\]
又 \(\widehat{\mathbf 1_S}(\{i\})=2^{-n}\sum_{\omega\in S}(-1)^{\omega_i}=b_i(S)/2^n\)，故
\[
2^{1-n}E_w(S)\ge \sum_{i=1}^n 4w_i\frac{b_i(S)^2}{2^{2n}},
\]
移项即得结论。证毕。
\end{proof}
\leanverified{paper\_fold\_hypercube\_weighted\_energy\_ellipsoid\_bias}

\begin{theorem}\label{thm:fold-hypercube-godel-ellipsoid-volume-energy-rigidity}
\leanverified{paper\_fold\_hypercube\_godel\_ellipsoid\_volume\_energy\_rigidity}
取互异素数 \(p_1,\dots,p_n\)，并定义对数权重能量
\[
E_{\log p}(S):=\sum_{i=1}^n (\log p_i)\,c_i(S),
\]
其中 \(c_i(S)\) 为第 \(i\) 坐标方向的无向边界数（命题 \ref{prop:fold-hypercube-godel-stokes-flux-bias}）。
对任意 \(S\subseteq\Omega_n\)，令
\(
r_i(S):=b_i(S)^2
\)
并构造轴对齐 Gödel 椭球
\[
\mathcal E_n(r(S))
:=\Bigl\{x\in\RR^n:\ \sum_{i=1}^n \frac{x_i^2}{p_i^{r_i(S)}}\le 1\Bigr\}.
\]
则其体积满足
\[
\log\operatorname{vol}\bigl(\mathcal E_n(r(S))\bigr)
=\log\kappa_n+\frac12\sum_{i=1}^n (\log p_i)\,r_i(S),
\]
并且成立严格上界
\[
\log\operatorname{vol}\bigl(\mathcal E_n(r(S))\bigr)
\le
\log\kappa_n+2^{n-2}\,E_{\log p}(S).
\]
若 \(S\) 非平凡（即 \(\varnothing\ne S\ne\Omega_n\)），则上式取等号当且仅当 \(S\) 为某一坐标的余维 \(1\) 面：
\[
S=\{\omega\in\Omega_n:\ \omega_j=\varepsilon\}
\quad\text{对某个 }j\in\{1,\dots,n\},\ \varepsilon\in\{0,1\}.
\]
\end{theorem}
\begin{proof}
体积公式来自轴对齐椭球的基本计算：
\(
\operatorname{vol}(\mathcal E_n(r))=\kappa_n\prod_{i=1}^n p_i^{r_i/2}
\)，取对数即得。
由命题 \ref{prop:fold-hypercube-weighted-energy-ellipsoid-bias} 取权重 \(w_i=\log p_i\) 得
\[
E_{\log p}(S)\ge 2^{1-n}\sum_{i=1}^n(\log p_i)\,b_i(S)^2
=2^{1-n}\sum_{i=1}^n(\log p_i)\,r_i(S),
\]
两边同乘 \(2^{n-2}\) 并代入体积公式即可得到上界。
\par
下面刻画等号情形。由命题 \ref{prop:fold-hypercube-weighted-energy-ellipsoid-bias} 的证明可见，
不等式来自谱展开中
\[
\sum_{\varnothing\ne I\subseteq[n]}4\Bigl(\sum_{i\in I}w_i\Bigr)|\widehat f(I)|^2
\ge
\sum_{i=1}^n 4w_i|\widehat f(\{i\})|^2,
\qquad f=\mathbf 1_S,
\]
故在 \(w_i=\log p_i>0\) 下取等号当且仅当 \(\widehat f(I)=0\) 对所有 \(|I|\ge 2\) 成立。
于是存在实数 \(a_0,a_1,\dots,a_n\) 使得
\[
f(\omega)=a_0+\sum_{i=1}^n a_i(-1)^{\omega_i}.
\]
由于 \(f(\omega)\in\{0,1\}\) 对一切 \(\omega\) 成立，若存在 \(i\ne j\) 使 \(a_i,a_j\ne 0\)，
则通过令 \((\omega_i,\omega_j)\) 取四种组合可使 \(f\) 取到 \(a_0\pm a_i\pm a_j\) 四个值，必有至少一值落在 \(\{0,1\}\) 之外，矛盾。
故至多一个 \(a_j\) 非零。
在 \(S\) 非平凡时必有某个 \(a_j\ne 0\)，并由二值约束可知
\(
f(\omega)=\frac12(1\pm (-1)^{\omega_j})
\)，等价于
\(
S=\{\omega:\ \omega_j=\varepsilon\}
\)。
反之，此类集合直接检验可使谱仅支撑在一次项上，从而取等号。证毕。
\end{proof}

\begin{theorem}[能量截止下的 Walsh--Stokes 全息重构与最优误差上界]\label{thm:fold-hypercube-walsh-stokes-energy-cutoff-reconstruction}
\leanverified{paper\_fold\_hypercube\_walsh\_stokes\_energy\_cutoff\_reconstruction}
在定理 \ref{thm:fold-hypercube-weighted-walsh-stokes} 的设定下，令 \(S\subseteq\Omega_n\)、\(f=\mathbf 1_S\)，并记
\[
\lambda_I:=4\sum_{i\in I}w_i,\qquad I\subseteq[n].
\]
对任意截止参数 \(\Lambda>0\)，令谱投影
\[
P_{\Lambda}f:=\sum_{\lambda_I\le \Lambda}\widehat f(I)\chi_I,
\qquad
R_{\Lambda}f:=f-P_{\Lambda}f.
\]
则在归一化 \(L^2\) 范数 \(\|g\|_2^2:=2^{-n}\sum_{\omega\in\Omega_n}g(\omega)^2\) 下成立误差界
\[
\boxed{\;\|R_{\Lambda}f\|_2^2\ \le\ \frac{\mathcal E_w(f)}{\Lambda}\ =\ \frac{2^{1-n}E_w(S)}{\Lambda}.\;}
\]
并且 \(P_\Lambda f\) 仅依赖于所有满足 \(\lambda_I\le \Lambda\) 的边界全息观测量
\(
\sum_{e\in\partial^{\to}S}\alpha_{I,w}(e)
\)
（定理 \ref{thm:fold-hypercube-weighted-walsh-stokes}），从而可由截止谱内的边界数据在 \(L^2\) 意义下重构 \(f\) 到误差不超过 \(2^{1-n}E_w(S)/\Lambda\)。
\end{theorem}
\begin{proof}
由 Parseval 恒等式，
\[
\|R_{\Lambda}f\|_2^2=\sum_{\lambda_I>\Lambda}|\widehat f(I)|^2.
\]
对每项用 \(\lambda_I>\Lambda\Rightarrow |\widehat f(I)|^2\le \Lambda^{-1}\lambda_I|\widehat f(I)|^2\) 得
\[
\|R_{\Lambda}f\|_2^2
\le
\frac{1}{\Lambda}\sum_{\lambda_I>\Lambda}\lambda_I|\widehat f(I)|^2
\le
\frac{1}{\Lambda}\sum_{I\ne\varnothing}\lambda_I|\widehat f(I)|^2.
\]
而 \(\sum_{I\ne\varnothing}\lambda_I|\widehat f(I)|^2=\mathcal E_w(f)=2^{1-n}E_w(S)\) 已在命题 \ref{prop:fold-hypercube-weighted-energy-ellipsoid-bias} 的谱展开口径中给出，故得到所示界。
最后由定理 \ref{thm:fold-hypercube-weighted-walsh-stokes} 的边界读出式 \(2^n\widehat f(I)=\sum_{e\in\partial^{\to}S}\alpha_{I,w}(e)\)，可知 \(P_\Lambda f\) 的全部系数由截止谱内的边界数据确定。证毕。
\end{proof}

\begin{definition}[斐波那契 Gödel 半径]\label{def:fold-hypercube-fibonacci-godel-radius}
令 \((F_k)_{k\ge 0}\) 为斐波那契数列（\(F_0=0,F_1=1\)），并取黄金分割比
\[
\varphi:=\frac{1+\sqrt5}{2}.
\]
对任意有限指标集 \(I\subseteq\NN\)，定义其斐波那契 Gödel 半径为
\[
\mathsf N_{\mathrm{Fib}}(I):=\sum_{i\in I}F_{i+1}.
\]
\end{definition}

\begin{proposition}[斐波那契 Gödel 半径的子集计数幂律]\label{prop:fold-hypercube-fibonacci-godel-radius-count}
\leanverified{paper\_fold\_hypercube\_fibonacci\_godel\_radius\_count}
令 \(\mathcal M_{\mathrm{Fib}}(R):=\{I\subseteq\NN:\ \mathsf N_{\mathrm{Fib}}(I)\le R\}\)（\(R\ge 1\)）。
则有
\[
\#\mathcal M_{\mathrm{Fib}}(R)=\Theta\!\left(R^{\frac{\log 2}{\log\varphi}}\right)
\qquad(R\to\infty).
\]
更精确地，若 \(k=k(R)\) 由 \(F_{k+1}\le R<F_{k+2}\) 唯一确定，则
\[
2^{k-2}\ \le\ \#\mathcal M_{\mathrm{Fib}}(R)\ \le\ 2^{k}.
\]
\end{proposition}
\begin{proof}
取 \(k\) 使 \(F_{k+1}\le R<F_{k+2}\)。
若 \(I\in\mathcal M_{\mathrm{Fib}}(R)\) 且 \(I\) 含有某个 \(i\ge k+1\)，则 \(\mathsf N_{\mathrm{Fib}}(I)\ge F_{i+1}\ge F_{k+2}>R\)，矛盾；
故 \(I\subseteq\{1,\dots,k\}\)，从而 \(\#\mathcal M_{\mathrm{Fib}}(R)\le 2^k\)。
\par
另一方面，对任意 \(J\subseteq\{1,\dots,k-2\}\) 有
\[
\mathsf N_{\mathrm{Fib}}(J)
\le
\sum_{i=1}^{k-2}F_{i+1}
=
\sum_{j=2}^{k-1}F_j
=
F_{k+1}-2
\le
R,
\]
其中使用恒等式 \(\sum_{j=1}^{m}F_j=F_{m+2}-1\)。
因此 \(\{1,\dots,k-2\}\) 的全部子集均属于 \(\mathcal M_{\mathrm{Fib}}(R)\)，给出下界 \(\#\mathcal M_{\mathrm{Fib}}(R)\ge 2^{k-2}\)。
\par
最后，由 \(F_{k+1}=\Theta(\varphi^{k})\) 得 \(k=\log_\varphi R+O(1)\)。
将 \(2^{k}\) 改写为 \(R^{\log_\varphi 2}\)（并用 \(\log_\varphi 2=\frac{\log 2}{\log\varphi}\)）即可得到 \(\#\mathcal M_{\mathrm{Fib}}(R)=\Theta\!\left(R^{\frac{\log 2}{\log\varphi}}\right)\)。
证毕。
\end{proof}

\begin{corollary}[斐波那契权重下的截止谱可见性与多项式维数律]\label{cor:fold-hypercube-fibonacci-godel-spectrum-polynomial}
\leanverified{paper\_fold\_hypercube\_fibonacci\_godel\_spectrum\_polynomial}
在定理 \ref{thm:fold-hypercube-walsh-stokes-energy-cutoff-reconstruction} 的设定下，取权重
\[
w_i:=F_{i+1},\qquad 1\le i\le n.
\]
则对每个 \(I\subseteq[n]\) 有
\[
\lambda_I=4\sum_{i\in I}w_i=4\,\mathsf N_{\mathrm{Fib}}(I).
\]
因此对任意 \(\varepsilon>0\)，若取
\[
\Lambda:=\frac{\mathcal E_w(f)}{\varepsilon^2},
\]
则由定理 \ref{thm:fold-hypercube-walsh-stokes-energy-cutoff-reconstruction} 可由所有满足 \(\mathsf N_{\mathrm{Fib}}(I)\le \Lambda/4\) 的边界全息观测确定 \(P_\Lambda f\)，并且
\[
\|f-P_\Lambda f\|_2^2=\|R_\Lambda f\|_2^2\le \varepsilon^2.
\]
当 \(n\) 足够大以容纳全部满足 \(\mathsf N_{\mathrm{Fib}}(I)\le R\) 的频率指标时，所需观测模态数满足
\[
\#\{I\subseteq[n]:\ \mathsf N_{\mathrm{Fib}}(I)\le R\}
=
\Theta\!\left(R^{\frac{\log 2}{\log\varphi}}\right)
\qquad(R\to\infty).
\]
\end{corollary}
\begin{proof}
由权重取值，\(\lambda_I=4\sum_{i\in I}F_{i+1}=4\,\mathsf N_{\mathrm{Fib}}(I)\) 立得截止条件等价性。
其余结论直接由定理 \ref{thm:fold-hypercube-walsh-stokes-energy-cutoff-reconstruction} 取 \(\Lambda=\mathcal E_w(f)/\varepsilon^2\) 与命题 \ref{prop:fold-hypercube-fibonacci-godel-radius-count} 给出。
\end{proof}

\begin{theorem}[斐波那契能量约束指示函数族的 \texorpdfstring{$\varepsilon$}{eps}-熵上界]\label{thm:fold-hypercube-fibonacci-eps-entropy-upper}
沿用推论 \ref{cor:fold-hypercube-fibonacci-godel-spectrum-polynomial} 的权重 \(w_i=F_{i+1}\) 与归一化 \(L^2\) 范数 \(\|\cdot\|_2\)。
对 \(E>0\) 定义能量受限类
\[
\mathcal F_E(n):=\{f=\mathbf 1_S:\ S\subseteq\Omega_n,\ \mathcal E_w(f)\le E\}\subset L^2(\Omega_n).
\]
记 \(N(\varepsilon,\mathcal F_E(n),\|\cdot\|_2)\) 为该类在距离 \(\|\cdot\|_2\) 下的 \(\varepsilon\)-覆盖数。
则存在与 \(n,E,\varepsilon\) 无关的常数 \(C_1,C_2>0\)，使对任意 \(0<\varepsilon<1\) 有
\[
\log N(\varepsilon,\mathcal F_E(n),\|\cdot\|_2)
\le
C_1\left(\frac{E}{\varepsilon^2}\right)^{\frac{\log 2}{\log\varphi}}
\log\!\left(C_2\frac{E}{\varepsilon}\right).
\]
\leanverified{paper\_fold\_hypercube\_fibonacci\_eps\_entropy\_upper}
\end{theorem}
\begin{proof}
对任意 \(f\in\mathcal F_E(n)\) 与任意 \(\Lambda>0\)，定理 \ref{thm:fold-hypercube-walsh-stokes-energy-cutoff-reconstruction} 给出
\[
\|f-P_\Lambda f\|_2^2\le \frac{\mathcal E_w(f)}{\Lambda}\le \frac{E}{\Lambda}.
\]
取 \(\Lambda:=4E/\varepsilon^2\) 得 \(\|f-P_\Lambda f\|_2\le \varepsilon/2\)。
因此构造 \(\varepsilon\)-网只需在有限维子空间
\[
V_\Lambda:=\mathrm{span}\{\chi_I:\ \lambda_I\le \Lambda\}
\]
上对 \(P_\Lambda f\) 作离散化。
\par
令
\[
N_\Lambda:=\dim V_\Lambda=\#\{I\subseteq[n]:\ \lambda_I\le \Lambda\}.
\]
由 \(\lambda_I=4\,\mathsf N_{\mathrm{Fib}}(I)\) 及命题 \ref{prop:fold-hypercube-fibonacci-godel-radius-count}（并注意 \(n\) 有限只会减小计数）可知
\[
N_\Lambda=\Theta\!\left(\left(\frac{\Lambda}{4}\right)^{\frac{\log 2}{\log\varphi}}\right)
=\Theta\!\left(\left(\frac{E}{\varepsilon^2}\right)^{\frac{\log 2}{\log\varphi}}\right).
\]
\par
写 \(P_\Lambda f=\sum_{\lambda_I\le \Lambda}\widehat f(I)\chi_I\)。
由于 Walsh 基在 \(\|\cdot\|_2\) 下正交归一，且 \(|\widehat f(I)|\le \|f\|_2\le 1\)，可对每个系数 \(\widehat f(I)\in[-1,1]\) 作步长
\[
\delta:=\frac{\varepsilon}{2\sqrt{N_\Lambda}}
\]
的均匀量化。
由量化误差的 \(\ell^2\) 估计，系数向量的误差满足
\(
\sqrt{N_\Lambda}\,\delta=\varepsilon/2
\)，从而在函数层面得到
\[
\|P_\Lambda f-\widetilde f\|_2\le \varepsilon/2.
\]
与 \(\|f-P_\Lambda f\|_2\le \varepsilon/2\) 合并即得 \(\|f-\widetilde f\|_2\le \varepsilon\)。
\par
每个系数区间 \([-1,1]\) 的网格点数至多 \(1+2/\delta\le 4/\delta\)（因 \(\varepsilon<1\) 且 \(N_\Lambda\ge 1\)）。
故覆盖数满足
\[
N(\varepsilon,\mathcal F_E(n),\|\cdot\|_2)
\le
\left(\frac{4}{\delta}\right)^{N_\Lambda}
=
\left(\frac{8\sqrt{N_\Lambda}}{\varepsilon}\right)^{N_\Lambda}.
\]
取对数并用 \(N_\Lambda=\Theta((E/\varepsilon^2)^{\frac{\log 2}{\log\varphi}})\) 吸收常数，得到所示上界形式。
\end{proof}

\begin{conjecture}[斐波那契谱指数的最坏情形刚性]\label{conj:fold-hypercube-fibonacci-spectral-exponent-rigidity}
在权重 \(w_i=F_{i+1}\) 下，记
\(
\mathcal M_{\mathrm{Fib}}(R)=\{I\subseteq\NN:\ \mathsf N_{\mathrm{Fib}}(I)\le R\}
\)
并由命题 \ref{prop:fold-hypercube-fibonacci-godel-radius-count} 知 \(\#\mathcal M_{\mathrm{Fib}}(R)\asymp R^{\frac{\log 2}{\log\varphi}}\)。
猜想存在绝对常数 \(c_0,c_1>0\)，使得对任意充分小的 \(\varepsilon\in(0,1)\) 与充分大的 \(E\)，存在维数 \(n\) 与某个集合 \(S\subseteq\Omega_n\) 满足 \(\mathcal E_w(\mathbf 1_S)\le E\)，并且对任意真子集
\[
\mathcal J\subsetneq \{I\subseteq[n]:\ \mathsf N_{\mathrm{Fib}}(I)\le c_0E/\varepsilon^2\}
\]
若 \(\#\mathcal J\le (1-c_1)\,\#\{I\subseteq[n]:\ \mathsf N_{\mathrm{Fib}}(I)\le c_0E/\varepsilon^2\}\)，则仅由 \(\{\widehat{\mathbf 1_S}(I)\}_{I\in\mathcal J}\)（等价：仅由相应边界全息观测）所构造的任意重构都不能保证在最坏情形下达到 \(L^2\) 误差 \(\le \varepsilon\)。
\par
等价地，指数 \(\frac{\log 2}{\log\varphi}\) 作为边界可见低能模态的增长阶在最坏情形下不可降低。
\end{conjecture}

\endinput


\endinput


% Quadratic multi-fiber Stokes--Plancherel area laws and audit corollaries.

\begin{theorem}[多纤维二次 Stokes--Plancherel 面积律与刚性]\label{thm:fold-coarsegraining-quadratic-stokes-area-rigidity}
\leanverified{paper\_fold\_coarsegraining\_quadratic\_stokes\_area\_rigidity}
在定义 \ref{def:hypercube-coarsegraining-area-gauge-volume} 的设定下，令 \(f:\Omega_n\twoheadrightarrow X\) 为满射粗粒化，并记纤维 \(S_x:=f^{-1}(x)\)（\(x\in X\)）。
对每个坐标 \(i\in\{1,\dots,n\}\)，定义纤维偏差
\[
b_i(S_x):=\card{S_x\cap\{\omega_i=0\}}-\card{S_x\cap\{\omega_i=1\}}
=\sum_{\omega\in S_x}(-1)^{\omega_i}\in\ZZ,
\]
并定义第 \(i\) 方向的切割边数
\[
c_i(f):=\card{\bigl\{\{\omega,\omega^{(i)}\}\in E_n:\ f(\omega)\neq f(\omega^{(i)})\bigr\}}.
\]
对任意全正权向量 \(w=(w_1,\dots,w_n)\in(0,\infty)^n\)，定义带权切割能量
\[
E_w(f):=\sum_{i=1}^n w_i\,c_i(f).
\]
则成立二次面积律
\[
\boxed{\;\sum_{x\in X}\sum_{i=1}^n w_i\,b_i(S_x)^2\ \le\ 2^n\,E_w(f).\;}
\]
特别地，当 \(w_i\equiv 1\) 时，
\[
\boxed{\;\sum_{x\in X}\sum_{i=1}^n b_i(S_x)^2\ \le\ 2^n\,A(f).\;}
\]

若 \(f\) 非平凡（非常值）且对某个全正权向量 \(w\) 上式取等号，则必有 \(|X|=2\)，并存在某个坐标 \(j\in\{1,\dots,n\}\) 使得在适当交换两类标签后
\[
S_{x_0}=\{\omega\in\Omega_n:\ \omega_j=0\},
\qquad
S_{x_1}=\{\omega\in\Omega_n:\ \omega_j=1\}.
\]
\end{theorem}
\begin{proof}
对任意纤维 \(S_x\subseteq\Omega_n\)，由命题 \ref{prop:fold-hypercube-weighted-energy-ellipsoid-bias}（取 \(S=S_x\)）得
\[
E_w(S_x):=\sum_{i=1}^n w_i\,c_i(S_x)\ \ge\ 2^{1-n}\sum_{i=1}^n w_i\,b_i(S_x)^2.
\]
对 \(x\in X\) 求和得
\begin{equation}\label{eq:fold-coarsegraining-quadratic-stokes-area-rigidity-step1}
\sum_{x\in X}E_w(S_x)\ \ge\ 2^{1-n}\sum_{x\in X}\sum_{i=1}^n w_i\,b_i(S_x)^2.
\end{equation}
另一方面，对固定坐标 \(i\)，每一条 \(i\)-方向切割边 \(\{\omega,\omega^{(i)}\}\)（即 \(f(\omega)\neq f(\omega^{(i)})\)）恰在端点所属的两条纤维边界中各计一次，故
\[
\sum_{x\in X}c_i(S_x)=2c_i(f).
\]
从而
\begin{equation}\label{eq:fold-coarsegraining-quadratic-stokes-area-rigidity-step2}
\sum_{x\in X}E_w(S_x)
=\sum_{i=1}^n w_i\sum_{x\in X}c_i(S_x)
=2\sum_{i=1}^n w_i\,c_i(f)
=2E_w(f).
\end{equation}
将 \eqref{eq:fold-coarsegraining-quadratic-stokes-area-rigidity-step2} 代回 \eqref{eq:fold-coarsegraining-quadratic-stokes-area-rigidity-step1} 即得二次面积律。
取 \(w_i\equiv 1\) 并注意 \(A(f)=\sum_i c_i(f)\) 得特例。
\par
现刻画等号情形。若 \(f\) 非常值且二次面积律取等号，则 \eqref{eq:fold-coarsegraining-quadratic-stokes-area-rigidity-step1} 亦取等号，从而每个纤维 \(S_x\) 都必须在命题 \ref{prop:fold-hypercube-weighted-energy-ellipsoid-bias} 的不等式中取等号。
由该命题证明中的 Walsh 谱展开可知：由于 \(w_i>0\)，取等号当且仅当 \(\widehat{\mathbf 1_{S_x}}(I)=0\) 对所有 \(|I|\ge 2\) 成立。
因此存在实数 \(\alpha,\beta_1,\dots,\beta_n\) 使得对一切 \(\omega\in\Omega_n\),
\[
\mathbf 1_{S_x}(\omega)=\alpha+\sum_{i=1}^n \beta_i\,(-1)^{\omega_i}.
\]
若存在 \(i\neq j\) 使 \(\beta_i,\beta_j\neq 0\)，则令 \((\omega_i,\omega_j)\) 取四种组合可使右端取到 \(\alpha\pm\beta_i\pm\beta_j\) 至少三个不同实值，与 \(\mathbf 1_{S_x}\in\{0,1\}\) 矛盾；故至多一个 \(\beta_j\) 非零。
在 \(S_x\) 非空且不为全体时，二值约束迫使 \(\alpha=\tfrac12\) 且 \(\beta_j=\pm\tfrac12\)，从而
\[
S_x=\{\omega\in\Omega_n:\ \omega_j=\varepsilon\}\qquad(\varepsilon\in\{0,1\}).
\]
由于 \(\{S_x\}_{x\in X}\) 两两不交且并为 \(\Omega_n\)，只能由一对互补半立方体组成，因而 \(|X|=2\)，并且两纤维对应同一坐标 \(j\) 的两侧，得到所述结论。证毕。
\end{proof}

\begin{corollary}[切割面积对 Gödel 外置码长平方和的控制]\label{cor:fold-coarsegraining-godel-length-square-area}
在定理 \ref{thm:fold-coarsegraining-quadratic-stokes-area-rigidity} 的设定下，取互异素数 \(p_1,\dots,p_n\)。
对每个纤维 \(S_x\) 定义外置码长
\[
L_{\boldsymbol p}(x):=\sum_{i=1}^n |b_i(S_x)|\log p_i.
\]
则成立估计
\[
\boxed{\;\sum_{x\in X}L_{\boldsymbol p}(x)^2\ \le\ 2^n\,A(f)\cdot\sum_{i=1}^n(\log p_i)^2.\;}
\]
\leanverified{paper\_fold\_coarsegraining\_godel\_length\_square\_area}
并且存在 \(x^\star\in X\) 满足
\[
\boxed{\;L_{\boldsymbol p}(x^\star)\ \ge\ \sqrt{\frac{2^nA(f)}{|X|}}\cdot\sqrt{\sum_{i=1}^n(\log p_i)^2}.\;}
\]
\end{corollary}
\begin{proof}
对任意 \(x\in X\)，由 Cauchy--Schwarz 不等式得
\[
L_{\boldsymbol p}(x)^2
=\Bigl(\sum_{i=1}^n |b_i(S_x)|\log p_i\Bigr)^2
\le
\Bigl(\sum_{i=1}^n b_i(S_x)^2\Bigr)\Bigl(\sum_{i=1}^n(\log p_i)^2\Bigr).
\]
对 \(x\) 求和并应用定理 \ref{thm:fold-coarsegraining-quadratic-stokes-area-rigidity} 在 \(w_i\equiv 1\) 的特例即可得到第一条不等式。
第二条由 \(\max_x L_{\boldsymbol p}(x)^2\ge |X|^{-1}\sum_x L_{\boldsymbol p}(x)^2\) 立得。证毕。
\end{proof}

\begin{theorem}[椭球约束下的 Gödel 码长对偶公式与素数--权重最优配对]\label{thm:fold-godel-ellipsoid-duality-optimal-pairing}
取全正权向量 \(w=(w_1,\dots,w_n)\in(0,\infty)^n\) 与半径 \(R>0\)。
对任意互异素数 \(p_1,\dots,p_n\) 与任意向量 \(b\in\RR^n\)，定义 Gödel 线性码长泛函
\[
\mathcal L_{\boldsymbol p}(b):=\sum_{i=1}^n |b_i|\log p_i.
\]
则：
\begin{enumerate}
  \item 精确对偶公式
  \[
  \boxed{\;
  \sup_{\sum_{i=1}^n w_i b_i^2\le R^2}\ \mathcal L_{\boldsymbol p}(b)
  \ =\ R\sqrt{\sum_{i=1}^n\frac{(\log p_i)^2}{w_i}}\; }.
  \]
  \item 在所有“选取 \(n\) 个互异素数并分配给坐标”的方案中，使右端最小的方案为：取最小的 \(n\) 个素数 \(q_1<\cdots<q_n\)，并按权重升序同序配对。
  更确切地，若 \(w_{(1)}\le\cdots\le w_{(n)}\) 为升序重排，则令
  \[
  p_{(k)}:=q_k\qquad(1\le k\le n)
  \]
  使得
  \(
  \sum_i(\log p_i)^2/w_i
  \)
  在所有方案中取到全局最小。
\end{enumerate}
\end{theorem}
\begin{proof}
\textbf{（对偶公式）}
由 Cauchy--Schwarz 不等式，
\[
\mathcal L_{\boldsymbol p}(b)
=\sum_{i=1}^n (\sqrt{w_i}|b_i|)\cdot\frac{\log p_i}{\sqrt{w_i}}
\le
\sqrt{\sum_{i=1}^n w_i b_i^2}\cdot \sqrt{\sum_{i=1}^n\frac{(\log p_i)^2}{w_i}}
\le
R\sqrt{\sum_{i=1}^n\frac{(\log p_i)^2}{w_i}}.
\]
取
\(
|b_i|=R\cdot\frac{\log p_i/w_i}{\sqrt{\sum_j (\log p_j)^2/w_j}}
\)
即可达上界，故得等号并得到精确值。
\par
\textbf{（最优配对）}
先固定一组互异素数集合 \(\{p_i\}\)，仅改变其在坐标间的排列。
令 \(a_i:=(\log p_i)^2\) 并按升序重排为 \(a_{(1)}\le\cdots\le a_{(n)}\)。
同时令 \(b_i:=1/w_i\)，若 \(w_{(1)}\le\cdots\le w_{(n)}\) 则 \(b_{(1)}\ge\cdots\ge b_{(n)}\)。
重排不等式给出：在所有排列 \(\sigma\) 中，
\(
\sum_{k=1}^n a_{(\sigma(k))}b_{(k)}
\)
的最小值由同序配对 \(a_{(k)}b_{(k)}\) 取得，即最大 \(a\) 配最小 \(b\)，等价于最大 \(\log p\) 配最大 \(w\)。
\par
再比较不同素数集合：由于 \(p\mapsto (\log p)^2\) 在 \(p>1\) 上严格递增，若将某个素数替换为更小的素数，则 \(\sum_i(\log p_i)^2/w_i\) 严格减小。
因此全局最小必由最小的 \(n\) 个素数给出，并结合上述排列极小化得到结论。证毕。
\end{proof}

\begin{corollary}[由二次偏差下界边界回路圆维]\label{cor:fold-coarsegraining-cycle-rank-from-quadratic-bias}
在定理 \ref{thm:fold-coarsegraining-quadratic-stokes-area-rigidity} 的设定下，记边界邻接多图为 \(\mathsf B_f\)（定义 \ref{def:hypercube-boundary-multigraph}），并记其连通分量数为 \(c(\mathsf B_f)\)。
则其圆维满足
\[
\boxed{\;
\cdim\bigl(H_1(\mathsf B_f;\ZZ)\bigr)
\ \ge\
2^{-n}\sum_{x\in X}\sum_{i=1}^n b_i(S_x)^2\ -\ |X|\ +\ c(\mathsf B_f).\;}
\]
\end{corollary}
\begin{proof}
由命题 \ref{prop:hypercube-boundary-multigraph-h1-cdim}，
\(
\cdim(H_1(\mathsf B_f;\ZZ))=A(f)-|X|+c(\mathsf B_f)
\)。
再用定理 \ref{thm:fold-coarsegraining-quadratic-stokes-area-rigidity} 的 \(w_i\equiv 1\) 特例给出
\(
A(f)\ge 2^{-n}\sum_{x,i} b_i(S_x)^2
\)，
代入即得。证毕。
\end{proof}

\begin{proposition}[面积对条件分布可辨识度的 \(L^2\) 控制]\label{prop:fold-coarsegraining-l2-distinguishability-area}
\leanverified{paper\_fold\_coarsegraining\_l2\_distinguishability\_area}
在定理 \ref{thm:fold-coarsegraining-quadratic-stokes-area-rigidity} 的设定下，令 \(W\) 为 \(\Omega_n\) 上的均匀随机变量，并记 \(Y:=f(W)\in X\)。
对每个坐标 \(i\in\{1,\dots,n\}\)，令
\[
P_i^{(0)}(x):=\PP(Y=x\mid W_i=0),\qquad
P_i^{(1)}(x):=\PP(Y=x\mid W_i=1),
\qquad x\in X.
\]
则对每个 \(i\) 与每个 \(x\in X\) 有恒等式
\[
P_i^{(0)}(x)-P_i^{(1)}(x)=\frac{b_i(S_x)}{2^{n-1}}.
\]
并且若对 \(u:X\to\RR\) 记 \(\|u\|_2^2:=\sum_{x\in X}u(x)^2\)，则成立面积上界
\[
\boxed{\;
\sum_{i=1}^n\bigl\|P_i^{(0)}-P_i^{(1)}\bigr\|_2^2
\ \le\ 2^{2-n}\,A(f).\;}
\]
\end{proposition}
\begin{proof}
由条件概率定义，
\[
P_i^{(0)}(x)=\frac{\card{S_x\cap\{\omega_i=0\}}}{2^{n-1}},
\qquad
P_i^{(1)}(x)=\frac{\card{S_x\cap\{\omega_i=1\}}}{2^{n-1}},
\]
相减即得第一式。
因此
\[
\bigl\|P_i^{(0)}-P_i^{(1)}\bigr\|_2^2
=\sum_{x\in X}\Bigl(\frac{b_i(S_x)}{2^{n-1}}\Bigr)^2
=2^{-2n+2}\sum_{x\in X}b_i(S_x)^2.
\]
对 \(i\) 求和并应用定理 \ref{thm:fold-coarsegraining-quadratic-stokes-area-rigidity} 的 \(w_i\equiv 1\) 特例即得
\[
\sum_{i=1}^n\bigl\|P_i^{(0)}-P_i^{(1)}\bigr\|_2^2
=2^{-2n+2}\sum_{x\in X}\sum_{i=1}^n b_i(S_x)^2
\le 2^{-2n+2}\cdot 2^n A(f)
=2^{2-n}A(f).
\]
证毕。
\end{proof}

\begin{corollary}[Gödel--Stokes 三角约束]\label{cor:fold-godel-stokes-triangle-constraint}
在推论 \ref{cor:fold-coarsegraining-godel-length-square-area} 的设定下，定义
\[
\mathsf{G}_{\boldsymbol p}(f):=\sum_{x\in X}L_{\boldsymbol p}(x)^2,
\qquad
\mathsf{D}(f):=\sum_{i=1}^n\bigl\|P_i^{(0)}-P_i^{(1)}\bigr\|_2^2,
\]
其中 \(P_i^{(0)},P_i^{(1)}\) 定义如命题 \ref{prop:fold-coarsegraining-l2-distinguishability-area}。
则成立无中间偏差变量的约束
\[
\boxed{\;
\mathsf{G}_{\boldsymbol p}(f)\ \le\ 2^n\Bigl(\sum_{i=1}^n(\log p_i)^2\Bigr)\,A(f),
\qquad
\mathsf{D}(f)\ \le\ 2^{2-n}\,A(f).\;}
\]
\end{corollary}
\begin{proof}
分别由推论 \ref{cor:fold-coarsegraining-godel-length-square-area} 与命题 \ref{prop:fold-coarsegraining-l2-distinguishability-area} 直接得到。证毕。
\end{proof}

\endinput


% Energy upper bounds for absolute-bias Gödel length.

\begin{theorem}[加权边界能量对 Gödel 码长的平方根上界]\label{thm:fold-hypercube-godel-length-energy-upperbound}
设 \(n\ge 1\)，\(\Omega_n=\{0,1\}^n\)，取 \(S\subseteq\Omega_n\)。
沿用命题 \ref{prop:fold-hypercube-godel-stokes-flux-bias} 的体积偏差 \(b_i(S)\) 与无向边界数 \(c_i(S)\)，并对权重 \(w_i>0\) 定义加权边界能量
\[
E_w(S):=\sum_{i=1}^n w_i\,c_i(S)
\]
（命题 \ref{prop:fold-hypercube-weighted-energy-ellipsoid-bias}）。
取互异素数 \(q_1,\dots,q_n\)，并定义
\[
N_S:=\prod_{i=1}^n q_i^{|b_i(S)|},\qquad
\Delta_S:=\frac12\log N_S.
\]
则成立估计
\[
\boxed{\;
\log N_S\ \le\ \sqrt{2^{\,n-1}E_w(S)}\;\sqrt{\sum_{i=1}^n \frac{(\log q_i)^2}{w_i}}\;}
\]
以及
\[
\boxed{\;
\Delta_S\ \le\ \frac12\sqrt{2^{\,n-1}E_w(S)}\;\sqrt{\sum_{i=1}^n \frac{(\log q_i)^2}{w_i}}\; }.
\]
特别地，若取 \(w_i:=\log q_i\)，则
\[
\log N_S\ \le\ \sqrt{2^{\,n-1}E_{\log q}(S)\sum_{i=1}^n \log q_i},
\qquad
\Delta_S\ \le\ \frac12\sqrt{2^{\,n-1}E_{\log q}(S)\sum_{i=1}^n \log q_i}.
\]
等价地，
\[
E_{\log q}(S)\ \ge\
2^{1-n}\cdot \frac{\bigl(\log N_S\bigr)^2}{\sum_{i=1}^n \log q_i}.
\]
\end{theorem}

\begin{proof}
由 Cauchy--Schwarz 不等式，
\[
\log N_S=\sum_{i=1}^n |b_i(S)|\log q_i
=\sum_{i=1}^n \bigl(\sqrt{w_i}|b_i(S)|\bigr)\cdot\frac{\log q_i}{\sqrt{w_i}}
\le
\Bigl(\sum_{i=1}^n w_i\,b_i(S)^2\Bigr)^{1/2}
\Bigl(\sum_{i=1}^n \frac{(\log q_i)^2}{w_i}\Bigr)^{1/2}.
\]
而命题 \ref{prop:fold-hypercube-weighted-energy-ellipsoid-bias} 给出
\(
E_w(S)\ge 2^{1-n}\sum_{i=1}^n w_i\,b_i(S)^2
\)，
故 \(\sum_i w_i b_i(S)^2\le 2^{n-1}E_w(S)\)。
合并即得第一条不等式；第二条由 \(\Delta_S=\tfrac12\log N_S\) 直接推出。
取 \(w_i=\log q_i\) 时，\(\sum_i (\log q_i)^2/w_i=\sum_i\log q_i\)，代入即得特例。证毕。
\end{proof}

\endinput



% Stokes--Plancherel boundary identity on the hypercube.

\begin{theorem}[Stokes--Plancherel 边界谱分解：割面积的完全通量公式]\label{thm:fold-hypercube-stokes-plancherel-boundary-energy}
\leanverified{paper\_fold\_hypercube\_stokes\_plancherel\_boundary\_energy}
设 \(n\ge 1\)，\(\Omega_n:=\{0,1\}^n\)。
令 \(S\subseteq\Omega_n\)，并记 \(f:=\mathbf 1_S\)。
对任意 \(a\in\Omega_n\) 定义 Walsh 字符
\[
\chi_a(\omega):=(-1)^{a\cdot \omega},
\qquad \omega\in\Omega_n,
\qquad |a|:=\#\{i:\ a_i=1\}.
\]
采用归一化 Fourier--Walsh 系数
\[
\widehat f(a):=2^{-n}\sum_{\omega\in\Omega_n} f(\omega)\chi_a(\omega).
\]
对每个 \(1\le i\le n\) 定义 Stokes 通量观测量
\[
J_{a,i}(S):=\sum_{\omega\in\Omega_n}\bigl(f(\omega)-f(\omega^{(i)})\bigr)\chi_a(\omega),
\]
其中 \(\omega^{(i)}\) 表示翻转第 \(i\) 位后的顶点。
则对每个 \(a\neq 0\) 成立：
\begin{enumerate}
  \item \label{it:fold-hypercube-stokes-readout} \textbf{（广义 Stokes 读出）}
  \[
  \widehat f(a)=\frac{2^{-(n+1)}}{|a|}\sum_{i:\ a_i=1}J_{a,i}(S).
  \]
  \item \label{it:fold-hypercube-plancherel-cut} \textbf{（边界 Plancherel 恒等式）}
  记 \(|\partial S|\) 为超立方体图上 \(S\) 的割边数，则
  \[
  |\partial S|
  =
  2^{n+1}\sum_{a\neq 0}|a|\,\widehat f(a)^2
  =
  2^{-(n+1)}\sum_{a\neq 0}\frac{1}{|a|}
  \Bigl(\sum_{i:\ a_i=1}J_{a,i}(S)\Bigr)^2.
  \]
\end{enumerate}
\end{theorem}
\begin{proof}
对任意 \(a\in\Omega_n\) 与任意 \(i\) 有
\[
\chi_a(\omega^{(i)})=
\begin{cases}
\chi_a(\omega),& a_i=0,\\
-\chi_a(\omega),& a_i=1.
\end{cases}
\]
因此对 \(a_i=1\)，变量替换 \(\omega\mapsto\omega^{(i)}\) 给出
\[
\sum_{\omega} f(\omega^{(i)})\chi_a(\omega)
=\sum_{\omega}f(\omega)\chi_a(\omega^{(i)})
=-\sum_{\omega}f(\omega)\chi_a(\omega),
\]
从而
\(
J_{a,i}(S)=2\sum_{\omega}f(\omega)\chi_a(\omega)=2^{n+1}\widehat f(a)
\)；
而当 \(a_i=0\) 时同理得 \(J_{a,i}(S)=0\)。
于是
\(
\sum_{i:\ a_i=1}J_{a,i}(S)=|a|\cdot 2^{n+1}\widehat f(a)
\)，得到 \eqref{it:fold-hypercube-stokes-readout}。
\par
为证明 \eqref{it:fold-hypercube-plancherel-cut}，注意对指示函数 \(f=\mathbf 1_S\) 有
\[
\sum_{\omega\in\Omega_n}\bigl(f(\omega)-f(\omega^{(i)})\bigr)^2=2\,c_i(S),
\]
其中 \(c_i(S)\) 为第 \(i\) 坐标方向的无向割边数；故
\[
|\partial S|=\sum_{i=1}^n c_i(S)=\frac12\sum_{i=1}^n\sum_{\omega\in\Omega_n}\bigl(f(\omega)-f(\omega^{(i)})\bigr)^2.
\]
另一方面，对差分算子 \(D_i f(\omega):=f(\omega)-f(\omega^{(i)})\) 用 Walsh--Parseval 恒等式并利用 \(a_i=1\) 时 \(D_i\chi_a=2\chi_a\)（否则为 \(0\)），得到
\[
2^{-n}\sum_{\omega}(D_i f(\omega))^2
=
\sum_{a\in\Omega_n} (2a_i)^2\,\widehat f(a)^2
=4\sum_{a:\ a_i=1}\widehat f(a)^2.
\]
乘以 \(2^n\) 后对 \(i\) 求和并代入上式，得
\[
|\partial S|
=\frac12\sum_{i=1}^n 2^n\cdot 4\sum_{a:\ a_i=1}\widehat f(a)^2
=2^{n+1}\sum_{a\neq 0}|a|\,\widehat f(a)^2.
\]
再将 \eqref{it:fold-hypercube-stokes-readout} 代入即得通量平方和表示。证毕。
\end{proof}

\begin{corollary}[平方偏差下界与低阶模式截断]\label{cor:fold-hypercube-stokes-plancherel-truncation}
\leanverified{paper\_fold\_hypercube\_stokes\_plancherel\_truncation}
在定理 \ref{thm:fold-hypercube-stokes-plancherel-boundary-energy} 的设定下：
\begin{enumerate}
  \item \textbf{（平方偏差下界）} 记位偏差 \(b_i(S):=\sum_{\omega\in S}(-1)^{\omega_i}\)，则
  \[
  |\partial S|\ \ge\ 2^{1-n}\sum_{i=1}^n b_i(S)^2.
  \]
  \item \textbf{（低阶模式截断与非负残差）} 对任意整数 \(k\ge 1\)，令
  \[
  L_k(S):=2^{-(n+1)}\sum_{\substack{a\neq 0\\ |a|\le k}}\frac{1}{|a|}
  \Bigl(\sum_{i:\ a_i=1}J_{a,i}(S)\Bigr)^2.
  \]
  则 \(|\partial S|\ge L_k(S)\)，并且差值具有严格非负分解
  \[
  |\partial S|-L_k(S)
  =
  2^{-(n+1)}\sum_{\substack{a\neq 0\\ |a|> k}}\frac{1}{|a|}
  \Bigl(\sum_{i:\ a_i=1}J_{a,i}(S)\Bigr)^2.
  \]
\end{enumerate}
\end{corollary}
\begin{proof}
第一条由定理 \ref{thm:fold-hypercube-stokes-plancherel-boundary-energy} 的通量平方和式保留 \(|a|=1\) 项并丢弃其余非负项得到。
注意当 \(a=e_i\) 时 \(\widehat f(e_i)=2^{-n}b_i(S)\)，并由定理 \ref{thm:fold-hypercube-stokes-plancherel-boundary-energy}\ref{it:fold-hypercube-stokes-readout} 得
\(
J_{e_i,i}(S)=2^{n+1}\widehat f(e_i)=2b_i(S)
\)；
代回即得系数 \(2^{1-n}\)。
第二条由同一定理的恒等式按 \(|a|\le k\) 与 \(|a|>k\) 分拆即可。证毕。
\end{proof}

\endinput

% Weighted Stokes--Sobolev energy identity and high-frequency tail bound.

\begin{theorem}[加权 Dirichlet 能量的 Walsh 谱恒等式与高频尾界]\label{thm:fold-hypercube-weighted-stokes-sobolev-energy-spectrum}
\leanverified{paper\_fold\_hypercube\_weighted\_stokes\_sobolev\_energy\_spectrum}
设 \(n\ge 1\)、\(\Omega_n:=\{0,1\}^n\)，并取权重向量 \(w=(w_1,\dots,w_n)\) 满足 \(w_i>0\)。
对任意函数 \(f:\Omega_n\to\CC\)，定义加权 Dirichlet 能量
\[
\mathcal E_w(f)
:=\frac12\sum_{\omega\in\Omega_n}\sum_{i=1}^n w_i\,\bigl|f(\omega^{(i)})-f(\omega)\bigr|^2.
\]
对任意 \(I\subseteq[n]\) 定义 Walsh 字符 \(\chi_I(\omega):=(-1)^{\sum_{i\in I}\omega_i}\) 并采用归一化 Fourier--Walsh 系数
\[
\widehat f(I):=2^{-n}\sum_{\omega\in\Omega_n} f(\omega)\chi_I(\omega).
\]
则成立精确谱分解
\[
\boxed{\;
\mathcal E_w(f)
=2^{n+1}\sum_{\varnothing\neq I\subseteq[n]}\Bigl(\sum_{i\in I}w_i\Bigr)\,|\widehat f(I)|^2.\;}
\]
特别地，若 \(f=\mathbf 1_S\) 为子集 \(S\subseteq\Omega_n\) 的指示函数，并记 \(c_i(S)\) 为第 \(i\) 坐标方向的无向割边数，则
\[
\mathcal E_w(\mathbf 1_S)=\sum_{i=1}^n w_i\,c_i(S).
\]
并且对任意阈值 \(\tau>0\) 有高频尾界
\[
\boxed{\;
\sum_{\substack{I\subseteq[n]\\ \sum_{i\in I}w_i\ge \tau}}|\widehat f(I)|^2
\ \le\ \frac{\mathcal E_w(f)}{2^{n+1}\tau}\; }.
\]
\end{theorem}
\begin{proof}
对加权 Laplace 算子取
\[
(\Delta_w g)(\omega):=\sum_{i=1}^n w_i\bigl(g(\omega)-g(\omega^{(i)})\bigr)
\]
（与定理 \ref{thm:fold-hypercube-weighted-walsh-stokes} 的记号一致）。
则由离散分部积分恒等式
\[
\langle f,\Delta_w f\rangle
=\frac12\sum_{\omega\in\Omega_n}\sum_{i=1}^n w_i\,|f(\omega)-f(\omega^{(i)})|^2
=\mathcal E_w(f),
\qquad
\langle f,g\rangle:=\sum_{\omega\in\Omega_n} f(\omega)\overline{g(\omega)}.
\]
另一方面，对任意 \(\varnothing\neq I\subseteq[n]\)，由
\(
\chi_I(\omega^{(i)})=\chi_I(\omega)
\)
当 \(i\notin I\)，以及
\(
\chi_I(\omega^{(i)})=-\chi_I(\omega)
\)
当 \(i\in I\)，得
\[
\Delta_w\chi_I=2\Bigl(\sum_{i\in I}w_i\Bigr)\chi_I,
\qquad
\langle \chi_I,\chi_I\rangle=2^n.
\]
将 \(f=\sum_{I\subseteq[n]}\widehat f(I)\chi_I\) 代入并用 Walsh 基正交性得到
\[
\mathcal E_w(f)=\langle f,\Delta_w f\rangle
=\sum_{\varnothing\neq I\subseteq[n]}|\widehat f(I)|^2\,\langle \chi_I,\Delta_w\chi_I\rangle
=2^{n+1}\sum_{\varnothing\neq I}\Bigl(\sum_{i\in I}w_i\Bigr)|\widehat f(I)|^2,
\]
即第一条恒等式。
\par
当 \(f=\mathbf 1_S\) 时，\(|\mathbf 1_S(\omega^{(i)})-\mathbf 1_S(\omega)|^2\) 等于 \(1\) 当且仅当 \(\omega\) 与 \(\omega^{(i)}\) 跨越割边；
对固定 \(i\) 的求和在每条无向割边上恰计两次，故
\(
\frac12\sum_{\omega}|\mathbf 1_S(\omega^{(i)})-\mathbf 1_S(\omega)|^2=c_i(S)
\)，从而 \(\mathcal E_w(\mathbf 1_S)=\sum_i w_i c_i(S)\)。
\par
最后，对满足 \(\sum_{i\in I}w_i\ge\tau\) 的 \(I\) 有
\(
|\widehat f(I)|^2\le \tau^{-1}\bigl(\sum_{i\in I}w_i\bigr)|\widehat f(I)|^2
\)；
对这类 \(I\) 求和并用谱恒等式即可得到尾界。证毕。
\end{proof}

\endinput

% Coordinate-resolved Stokes torsor structure on boundary multigraphs.

\begin{corollary}[坐标分辨的 Stokes 扭量：向量散度约束的同调与上同调口径]\label{cor:fold-boundary-stokes-torsor-h1-h1}
\leanverified{paper\_fold\_boundary\_stokes\_torsor\_h1\_h1}
沿用定义 \ref{def:hypercube-boundary-multigraph} 的边界邻接多图 \(\mathsf B_f\)，并按命题 \ref{prop:fold-hypercube-godel-hodge-torsor} 记每个坐标的纤维偏差向量 \(b^{(i)}\in C_0(\mathsf B_f;\ZZ)\) 与整数流解集
\[
\mathcal S_i=\{F\in C_1(\mathsf B_f;\ZZ):\ \partial F=b^{(i)}\}.
\]
将各坐标并置为向量值数据
\[
b:=(b^{(1)},\dots,b^{(n)})\in C_0(\mathsf B_f;\ZZ^n),
\qquad
\mathcal S_b:=\{F\in C_1(\mathsf B_f;\ZZ^n):\ \partial F=b\}.
\]
则若 \(\mathcal S_b\neq\varnothing\)，解集为仿射格
\[
\mathcal S_b=F^{(0)}+\ker(\partial)\otimes_{\ZZ}\ZZ^n,
\]
从而在平移群 \(H_1(\mathsf B_f;\ZZ)^n\cong \ker(\partial)\otimes\ZZ^n\) 作用下构成扭量。
由于 \(\mathsf B_f\) 为一维 CW 复形且 \(H_1(\mathsf B_f;\ZZ)\) 自由阿贝尔，亦有自然同构
\[
H_1(\mathsf B_f;\ZZ)^n
\ \cong\
H^1(\mathsf B_f;\ZZ^n),
\]
从而上述扭量结构可等价地表述为 \(H^1(\mathsf B_f;\ZZ^n)\) 扭量。
\end{corollary}
\begin{proof}
命题 \ref{prop:fold-hypercube-godel-hodge-torsor} 对每个坐标给出非空时的扭量分解 \(\mathcal S_i=F^{(i)}+\ker(\partial)\)。
将 \(n\) 个坐标取直积即得
\(
\mathcal S_b=\prod_i \mathcal S_i = F^{(0)}+\ker(\partial)\otimes\ZZ^n
\)
（取 \(F^{(0)}:=(F^{(1)},\dots,F^{(n)})\)）。
又对图的链复形有 \(H_1(\mathsf B_f;\ZZ)\cong \ker(\partial)\)，且 \(H_1(\mathsf B_f;\ZZ)\) 自由使得
\(
H^1(\mathsf B_f;\ZZ^n)\cong \Hom(H_1(\mathsf B_f;\ZZ),\ZZ^n)\cong H_1(\mathsf B_f;\ZZ)^n
\)，
从而得到上同调表述。证毕。
\end{proof}

\endinput

% Vector-valued Hodge--Stokes orthogonal decomposition on boundary graphs.

\begin{corollary}[Hodge--Stokes 正交分解与最小能量原理（向量值版本）]\label{cor:fold-boundary-hodge-stokes-orthogonal-vector}
\leanverified{paper\_fold\_boundary\_hodge\_stokes\_orthogonal\_vector}
令 \(G=(V,E)\) 为有限连通无向多重图，固定一组取向并把 \(C_1(G;\RR^n)\) 与 \(C_0(G;\RR^n)\) 视作有限维 Hilbert 空间，内积取为
\[
\langle g,h\rangle:=\sum_{e\in E} g(e)\cdot h(e).
\]
记边界算子 \(\partial:C_1(G;\RR^n)\to C_0(G;\RR^n)\)，并令 \(\partial^\ast\) 为其伴随，图 Laplacian 为 \(L:=\partial\partial^\ast\)。
对任意散度数据 \(b\in C_0(G;\RR^n)\) 满足守恒条件 \(\sum_{v\in V} b(v)=0\)，考虑约束极小化问题
\[
\min\{\|g\|_2^2:\ g\in C_1(G;\RR^n),\ \partial g=b\}.
\]
则存在唯一最小能量解 \(g_{\min}\)，并且：
\begin{enumerate}
  \item \(g_{\min}\in \mathrm{im}(\partial^\ast)\)，从而 \(g_{\min}=\partial^\ast\varphi\) 且 \(L\varphi=b\)，其中 \(\varphi\) 在去除常数自由度后唯一。
  \item 对任意满足 \(\partial g=b\) 的 \(g\)，存在唯一分解
  \[
  g=g_{\min}+h,
  \qquad h\in \ker(\partial),
  \]
  并满足严格毕达哥拉斯恒等式
  \[
  \|g\|_2^2=\|g_{\min}\|_2^2+\|h\|_2^2.
  \]
\end{enumerate}
当取 \(G=\mathsf B_f\) 且 \(b\) 由纤维偏差向量值数据给出时，上述 \(h\) 的同调类即为边界 Stokes 扭量自由度的连续化刻画。
\end{corollary}
\begin{proof}
约束集合 \(\{g:\ \partial g=b\}\) 为仿射闭子空间，且与 \(\ker(\partial)\) 平行。
在 Hilbert 空间中，最小范数元为该仿射子空间到原点的正交投影，故满足
\(
g_{\min}\perp \ker(\partial)
\)
且 \(\partial g_{\min}=b\)。
又有限维情形有 \((\ker(\partial))^\perp=\mathrm{im}(\partial^\ast)\)，故存在 \(\varphi\) 使 \(g_{\min}=\partial^\ast\varphi\)，代入约束得 \(L\varphi=b\)。
由于 \(G\) 连通，\(\ker(L)=\langle \mathbf 1\rangle\otimes\RR^n\)，故去除常数自由度后 \(\varphi\) 唯一。
任取 \(g\) 满足 \(\partial g=b\)，则 \(h:=g-g_{\min}\in\ker(\partial)\)，并由 \(g_{\min}\perp h\) 得能量正交分解。证毕。
\end{proof}

\endinput

% Naturality under refinement and Rayleigh-type monotonicity for Stokes energies.

\begin{proposition}[粗粒化细化下的 Stokes 自然性与 Rayleigh--Stokes 单调性]\label{prop:fold-coarsegraining-naturality-rayleigh-stokes}
\leanverified{paper\_fold\_coarsegraining\_naturality\_rayleigh\_stokes}
设 \(n\ge 1\) 且 \(\Omega_n=\{0,1\}^n\)。
给定两级满射粗粒化 \(\tilde f:\Omega_n\twoheadrightarrow \tilde X\) 与 \(f:\Omega_n\twoheadrightarrow X\)，并存在满射 \(\rho:\tilde X\twoheadrightarrow X\) 使 \(f=\rho\circ \tilde f\)。
记 \(\mathsf B_{\tilde f}\)、\(\mathsf B_f\) 为相应边界邻接多图（定义 \ref{def:hypercube-boundary-multigraph}）。
\begin{enumerate}
  \item \textbf{（坐标净通量的自然性）}
  令 \(F^{(i)}_{\tilde f}\in C_1(\mathsf B_{\tilde f};\ZZ)\) 与 \(F^{(i)}_{f}\in C_1(\mathsf B_{f};\ZZ)\) 为命题 \ref{prop:fold-hypercube-godel-hodge-torsor} 给出的第 \(i\) 坐标自然整数流。
  定义推前算子 \(\rho_\#:C_1(\mathsf B_{\tilde f};\ZZ)\to C_1(\mathsf B_f;\ZZ)\) 为：对任意有向边 \(u\to v\)（\(u,v\in X\)）取
  \[
  (\rho_\#\Phi)(u\to v):=\sum_{\substack{\tilde u\in\rho^{-1}(u)\\ \tilde v\in\rho^{-1}(v)}}\Phi(\tilde u\to \tilde v).
  \]
  则对每个 \(i\) 有 \(\rho_\#F^{(i)}_{\tilde f}=F^{(i)}_{f}\)，并且边界算子满足交换律 \(\partial(\rho_\#\Phi)=\rho_\#(\partial\Phi)\)。
  \item \textbf{（Rayleigh--Stokes 单调性）}
  在 \(\mathsf B_{\tilde f}\) 的每条边 \(\tilde e\) 上取导纳 \(\tilde c_{\tilde e}>0\)，并对粗图边 \(e\) 定义推前导纳
  \[
  c_e:=\sum_{\tilde e\mapsto e}\tilde c_{\tilde e}.
  \]
  在 \(C_1(\mathsf B_f;\RR^n)\) 上定义能量
  \[
  \mathcal E_c(g):=\sum_{e\in E(\mathsf B_f)}\frac{\|g(e)\|^2}{c_e},
  \]
  并对守恒散度数据 \(b\in C_0(\mathsf B_f;\RR^n)\) 定义最小能量
  \[
  \mathfrak E_f(b):=\min\{\mathcal E_c(g):\ g\in C_1(\mathsf B_f;\RR^n),\ \partial g=b\}.
  \]
  在细图上同理定义 \(\mathfrak E_{\tilde f}(\tilde b)\)。
  令细化分配集合
  \[
  \mathcal B_{\tilde f}(b):=\{\tilde b\in C_0(\mathsf B_{\tilde f};\RR^n):\ \rho_\#\tilde b=b\}.
  \]
  则成立不等式
  \[
  \inf_{\tilde b\in\mathcal B_{\tilde f}(b)}\ \mathfrak E_{\tilde f}(\tilde b)\ \ge\ \mathfrak E_f(b).
  \]
\end{enumerate}
\end{proposition}
\begin{proof}
\textbf{（自然性）}
由定义，粗纤维按细纤维分解；超立方体第 \(i\) 坐标割边推前到边界邻接多图时，跨越粗纤维的割边是跨越相应细纤维割边的不交并，且取向保持。
因此坐标整数流在聚合下逐边相加，得到 \(\rho_\#F^{(i)}_{\tilde f}=F^{(i)}_{f}\)。
边界算子交换律为定义的直接重排。
\par
\textbf{（单调性）}
对任意 \(\tilde b\in\mathcal B_{\tilde f}(b)\)，取任意 \(\tilde g\in C_1(\mathsf B_{\tilde f};\RR^n)\) 满足 \(\partial\tilde g=\tilde b\)，并令 \(g:=\rho_\#\tilde g\)。
由 \(\partial(\rho_\#\cdot)=\rho_\#(\partial\cdot)\) 得 \(\partial g=b\)。
对每条粗边 \(e\) 记提升边集 \(\{\tilde e:\tilde e\mapsto e\}\)，则向量值 Cauchy--Schwarz 给出
\[
\frac{\|g(e)\|^2}{c_e}
=\frac{\bigl\|\sum_{\tilde e\mapsto e}\tilde g(\tilde e)\bigr\|^2}{\sum_{\tilde e\mapsto e}\tilde c_{\tilde e}}
\le
\sum_{\tilde e\mapsto e}\frac{\|\tilde g(\tilde e)\|^2}{\tilde c_{\tilde e}}.
\]
对全部粗边求和得 \(\mathcal E_c(g)\le \mathcal E_{\tilde c}(\tilde g)\)，从而
\[
\mathfrak E_f(b)\le \mathcal E_c(g)\le \mathcal E_{\tilde c}(\tilde g).
\]
对满足 \(\partial\tilde g=\tilde b\) 的 \(\tilde g\) 取极小，再对 \(\tilde b\in\mathcal B_{\tilde f}(b)\) 取下确界，即得结论。证毕。
\end{proof}

\endinput

% Boundary modularity conjecture linking Ihara zeta and Lee--Yang elliptic factors.

\begin{conjecture}[边界 Stokes 扭量的模性兼容性]\label{conj:fold-boundary-stokes-langlands-modularity}
设 \(\{f_m:\Omega_n\twoheadrightarrow X_m\}_{m\ge 1}\) 为由递归细化/时间尺度机制诱导的一族粗粒化，并记其边界邻接多图为 \(\mathsf B_{f_m}\)。
设与 Lee--Yang 完成结构相容的椭圆曲线族 \(\{E_m\}\) 控制相应零集的代数几何接口（例如正文中关于椭圆结构与模性链的口径）。
则存在一组谱归一化与局部因子选取，使得对每个 \(m\)：
\begin{enumerate}
  \item \(\mathsf B_{f_m}\) 的 Ihara 型 \(\zeta\) 函数的非平凡因子与 \(E_m\) 的 Hasse--Weil \(L\)-函数在素数处的局部因子在特征多项式意义下相容；
  \item 结论 \ref{cor:fold-boundary-stokes-torsor-h1-h1} 所刻画的 Stokes 扭量自由度在模 \(\ell\) 约化后，给出与 \(E_m\) 的 \(\ell\)-进表示的某一交换化截面相匹配的可观测子商。
\end{enumerate}
\end{conjecture}

\endinput


\endinput
